\documentclass[10pt]{amsart}
\usepackage[a4paper,margin=20mm]{geometry}
\emergencystretch=3em
\usepackage{amsmath,amssymb,amsthm,mathtools,mathrsfs,enumitem,hyperref}
\newtheorem{theo}{Theorem}[section]
\newtheorem{prop}[theo]{Proposition}
\newtheorem{lemm}[theo]{Lemma}
\newtheorem{coro}[theo]{Corollary}
\theoremstyle{definition}\newtheorem{defi}[theo]{Definition}
\theoremstyle{remark}\newtheorem{rema}[theo]{Remark}
\newcommand{\firstname}[1]{#1}\newcommand{\lastname}[1]{#1}
\newcommand{\eg}{e.g.\ }\newcommand{\ie}{i.e.\ }\newcommand{\resp}{resp.\ }\newcommand{\cf}{cf.\ }

\DeclarePairedDelimiter\abs{\lvert}{\rvert} 

\newcommand{\Fg}{\mathfrak{g}}
\newcommand{\Fh}{\mathfrak{h}}

\newcommand{\BC}{\mathbb{C}}
\newcommand{\BR}{\mathbb{R}}
\newcommand{\BQ}{\mathbb{Q}}
\newcommand{\BZ}{\mathbb{Z}}
\newcommand\bm[1]{\boldsymbol{#1}}
\newcommand\bmrho[1]{\boldsymbol{#1}}




\newcommand{\U}{\mathbf{U}}
\newcommand{\si}{\frac{\infty}{2}}


\DeclareMathOperator{\af}{af}
\DeclareMathOperator{\wt}{wt}
\DeclareMathOperator{\LT}{LT}
\DeclareMathOperator{\CST}{CST}
\DeclareMathOperator{\Hom}{Hom}
\DeclareMathOperator{\Bij}{Bij}
\DeclareMathOperator{\Par}{Par}
\DeclareMathOperator{\SiB}{SiB}



\usepackage{ytableau}
\ytableausetup{centertableaux,mathmode,smalltableaux}
\def\young(#1){\ytableaushort{#1}}
\def\yng(#1){\ydiagram{#1}}



%%%%%%%%% a placer avant \begin{document}


%%%%%%%%%%%%%%%%%%%%%%%%%%%%%%%%%%%%%
\newcommand*{\mk}{\mkern -1mu}
\newcommand*{\Mk}{\mkern -2mu}
\newcommand*{\mK}{\mkern 1mu}
\newcommand*{\MK}{\mkern 2mu}

\hypersetup{urlcolor=purple, linkcolor=blue, citecolor=red}


\newcommand*{\romanenumi}{\renewcommand*{\theenumi}{\roman{enumi}}}
\newcommand*{\Romanenumi}{\renewcommand*{\theenumi}{\Roman{enumi}}}
\newcommand*{\alphenumi}{\renewcommand*{\theenumi}{\alph{enumi}}}
\newcommand*{\Alphenumi}{\renewcommand*{\theenumi}{\Alph{enumi}}}
%%%%%%%%%%%%%%%%%%%%%%%%%%%%%%%%%%



%%%%%%%%%%%%%%%%%%%%%%%%%%%%%%%%%%%%%%%%%%%
%%%%% Auteur

\author{\firstname{Motohiro} \lastname{Ishii}}

\address{Department of Mathematics\\
Cooperative Faculty of Education\\
Gunma University\\ 
Maebashi\\
Gunma\\
371-8510\\
Japan}

\email{m.ishii@gunma-u.ac.jp}

\thanks{This work was supported by JSPS KAKENHI Grant-in-Aid for Young Scientists (B) 16K17577, Japan.}


%%%%% Sujet

\subjclass{}

\keywords{Semi-infinite Young tableau, semi-infinite Lakshmibai--Seshadri path, semi-infinite Bruhat order, affine Weyl group, quantum affine algebra, extremal weight module, crystal basis.}



%%%%% Gestion







%%%%% Titre et résumé
\title
[Semi-infinite Young tableaux]
{Semi-infinite Young tableaux and standard monomial theory for semi-infinite Lakshmibai--Seshadri paths}

\begin{document}

\begin{abstract}
We introduce semi-infinite Young tableaux, 
and show that these tableaux give a 
combinatorial model for the crystal basis 
of a level-zero extremal weight module 
over the quantized universal enveloping algebra of untwisted affine type $A$. 
The definition and characterization of these tableaux are based on 
standard monomial theory for semi-infinite Lakshmibai--Seshadri paths 
and a tableau criterion for the semi-infinite Bruhat order 
on affine Weyl groups of type $A$, 
which are also proved in this paper. 
\end{abstract}


\maketitle
\section{Introduction}

The aim of this paper is to introduce semi-infinite Young tableaux
(see Definition~\ref{def:SiYT}~\ref{Defi4.2_2}). 
These tableaux give a new combinatorial model for 
the crystal basis of a level-zero extremal weight module 
(see \S~\ref{Subsection:extremal})
over the quantized universal enveloping algebra of untwisted affine type $A$. 
In order to accomplish our purpose, we investigate 
\begin{enumerate}[label=(\roman*)]
\item
a characterization of 
the image of the strict embedding 
$\Phi_{\lambda \mid q=0}^{\LT }$ (see \S~\ref{Subsection:LT})
of crystals in terms of the semi-infinite Bruhat order 
(see Theorem~\ref{thm:SMT}), and 
\item
a tableau criterion for the semi-infinite Bruhat order on 
affine Weyl groups of type $A$ in Grassmannian cases (see Theorem~\ref{thm:tab-cri}).
\end{enumerate}
Note that the image of 
$\Phi_{\lambda \mid q=0}^{\LT }$
is an isomorphic image of the crystal basis of 
the extremal weight module of extremal weight $\lambda$ (of level-zero) 
into the tensor product of crystal bases of extremal weight modules 
associated with level-zero fundamental weights. 

Various generalizations and variations of (semi-standard) Young tableaux 
are concerned in many areas such as 
algebraic combinatorics, representation theory, algebraic geometry, and so forth. 
In particular, 
Littelmann (\cite{L94,L95}) introduced the Lakshmibai--Seshadri paths 
for all symmetrizable Kac--Moody root data, 
which can be thought of as a type-free generalization of Young tableaux. 
Soon after, Joseph (\cite{J}) and Kashiwara (\cite{K96}) independently proved that, 
for a dominant integral weight $\Lambda$ of a symmetrizable Kac--Moody Lie algebra $\Fg$, 
the set of Lakshmibai--Seshadri paths of shape $\Lambda$ 
equipped with Littelmann's root operators 
is isomorphic, as a $\Fg$-crystal, to the crystal basis of 
the integrable (irreducible) highest weight module of highest weight $\Lambda$ 
over the quantized universal enveloping algebra associated with $\Fg$. 
In view of Kashiwara's crystal (basis) theory, 
further generalizations and variants of Littelmann's path model are investigated; \eg generalized 
Lakshmibai--Seshadri paths for Borcherds--Kac--Moody root data (\cite{I,JL}), 
quantum Lakshmibai--Seshadri paths and 
semi-infinite Lakshmibai--Seshadri paths for untwisted affine root data
(\cite{INS16,LNSSS16}). 
Among these general theories, 
it should be emphasized that the original Young tableaux have especially nice combinatorial structures 
(see for instance~\cite{F,Kw})
due to the fact that every fundamental representation of finite type $A$ is minuscule; 
namely, the Weyl group acts transitively 
on the crystal basis of any fundamental representation in the case of finite type $A$. 
Similarly, in the case of untwisted affine type $A$, 
every extremal weight module associated with a 
level-zero fundamental weight is minuscule (see Proposition~\ref{prop:aff-varpi_i}~\ref{Prop4.3_2}); 
in this case, the affine Weyl group acts transitively on the crystal basis. 
Therefore, it is natural to try to find a tableau model 
for crystal bases of level-zero extremal weight modules 
in the case of untwisted affine type $A$. 

Let us give an explanation of our strategy to introduce semi-infinite Young tableaux. 

For this purpose, 
we first briefly sketch a standard monomial theoretic characterization 
of (ordinary) Young tableaux 
in terms of crystal basis theory as follows: 
let 
$U_q(\mathfrak{sl}_n(\BC))$ 
be the quantized universal enveloping algebra of type $A_{n-1}$ (see~\cite[\S~4.8 TABLE Fin]{Kac}).
Let $\varpi_i$, $1 \leq i \leq n-1$, be the $i$-th fundamental weight; 
we identify a dominant integral weight 
$\lambda = \sum_{i=1}^{n-1} m_i \varpi_i$, 
$m_i \in \BZ_{\geq 0}$, $1 \leq i \leq n-1$, 
with the Young diagram such that 
the number of the columns of length $i$ is $m_i$ for $1 \leq i \leq n-1$
(see Remark~\ref{rem:Young}). 
For a dominant integral weight 
$\lambda$, 
let 
$L(\lambda)$ 
be the irreducible finite-dimensional highest weight 
$U_q(\mathfrak{sl}_n(\BC))$-module of highest weight 
$\lambda$, 
and let 
$B(\lambda)$ 
be the crystal basis of 
$L(\lambda)$. 
It follows that 
$B(\varpi_i)$ 
is parametrized by the set 
$\CST (\varpi_i)$ 
of column-strict tableaux of shape 
$\varpi_i$
with entries in 
$\{ 1,2,\ldots ,n \}$
(see for instance~\cite[Proposition~3.3.1~(i)]{KN}). 
We have an injective homomorphism 
\begin{align} \label{eq:emb-hwm}
L(\lambda) \longrightarrow \bigotimes_{i=1}^{n-1} L(\varpi_i)^{\otimes m_i}
\end{align}
of $U_q(\mathfrak{sl}_n(\BC))$-modules 
sending a highest weight vector to the tensor product of highest weight vectors. 
Further, this homomorphism induces a strict embedding 
\begin{align} \label{eq:emb-sst}
B(\lambda) \longrightarrow \bigotimes_{i = 1}^{n-1} B(\varpi_i)^{\otimes m_i} 
\, \cong \, 
\prod_{i=1}^{n-1} \CST (\varpi_i)^{m_i}
\, \cong \, 
\CST (\lambda) 
\end{align}
of $U_q(\mathfrak{sl}_n(\BC))$-crystals, 
where 
$\CST (\lambda)$ 
denotes the set of column-strict tableaux of shape $\lambda$ with entries in $\{ 1,2,\ldots ,n \}$. 
Since the symmetric group of degree $n$ acts transitively on 
$\CST (\varpi_i)$, 
each element in 
$\CST (\lambda)$ 
is labeled by a tuple of 
$N := \sum_{i=1}^{n-1} m_i$ 
cosets in the symmetric group; 
the symmetric group of degree $n$ will be 
viewed as the Weyl group of type $A_{n-1}$. 
Let $\mathbb{T} \in \CST (\lambda)$, 
and assume that $\mathbb{T}$ corresponds to the tuple 
$(\bar{w}_1 , \bar{w}_2 , \ldots , \bar{w}_N)$
of $N$ cosets in the symmetric group. 
Then 
$\mathbb{T}$ 
is in the image of the strict embedding~\eqref{eq:emb-sst} 
if and only if there exist coset representatives 
$w_{\nu} \in \bar{w}_{\nu}$, $1 \leq \nu \leq N$, 
such that
$w_1 \succeq w_2 \succeq \cdots \succeq w_N$ 
in the Bruhat order $\succeq$ on the symmetric group
(see~\cite[Theorem~10.1]{L96}). 
In consequence, a column-strict tableau satisfying this condition 
is just a Young tableau, and vise versa
(\cite[Theorem~2.6.3 (Tableau Criterion)]{BB}; 
see also~\cite[Theorem~3.4.2~(i)]{KN}). 
This gives an isomorphism of 
$U_q(\mathfrak{sl}_n(\BC))$-crystals 
between 
$B(\lambda)$ 
and the set of Young tableaux of shape $\lambda$. 
For an explicit description of the $U_q(\mathfrak{sl}_n(\BC))$-crystal structure on
the set of Young tableaux, see~\cite[Theorem~3.4.2~(ii)]{KN}.

In the case of untwisted affine type $A$, 
our basic idea is to use 
the semi-infinite Bruhat order on the affine Weyl group 
in place of the Bruhat order on the symmetric group. 
Let 
$\U$ 
be the quantized universal enveloping algebra of (arbitrary) untwisted affine type, 
and let 
$I_{\af} = \{ 0 \} \sqcup I$
be the index set for the simple roots.
By abuse of notation, we use the same symbol 
$\varpi_i$, $i \in I$, 
for the $i$-th level-zero fundamental weight of $\U$. 
For 
$\lambda = \sum_{i \in I} m_i \varpi_i$, 
with 
$m_i \in \BZ_{\geq 0}$, $i \in I$, 
let 
$V(\lambda)$ 
be the extremal weight $\U$-module 
of extremal weight $\lambda$, 
and let 
$\mathcal{B}(\lambda)$ 
be the crystal basis of 
$V(\lambda)$ 
(see for instance~\cite[\S~3.1]{K02}). 
Similarly to~\eqref{eq:emb-hwm}, 
we have a canonical homomorphism 
\begin{align}
\Phi_{\lambda} : 
V(\lambda) 
\longrightarrow 
\bigotimes_{i \in I} V(\varpi_i)^{\otimes m_i}
\end{align}
of $\U$-modules sending an extremal weight vector to the tensor product of extremal weight vectors. 
The main difficulty in carrying out the argument similar to the above is that the associated map 
$\Phi_{\lambda \mid q=0}$ 
at 
$q = 0$ 
of 
$\Phi_{\lambda}$ 
does not necessarily induce a morphism 
$\mathcal{B}(\lambda) \rightarrow \bigotimes_{i \in I} \mathcal{B}(\varpi_i)^{\otimes m_i}$
of $\U$-crystals. 
In fact, each 
$\Phi_{\lambda \mid q=0}(b)$, 
$b \in \mathcal{B}(\lambda)$, 
is a linear combination of crystal basis elements 
whose terms are in one-to-one correspondence with 
the terms of a product of some Schur polynomials 
(see~\cite[\S~4.2]{BN04}). 
To overcome this difficulty, 
by taking the ``leading term'' 
$\Phi_{\lambda \mid q=0}^{\LT }(b)$ 
of 
$\Phi_{\lambda \mid q=0}(b)$
(see Remark~\ref{rem:LT}), 
we introduce a strict embedding 
\begin{align} \label{eq:intro-LT}
\Phi_{\lambda \mid q=0}^{\LT } : 
\mathcal{B}(\lambda) 
\longrightarrow 
\bigotimes_{i \in I} \mathcal{B}(\varpi_i)^{\otimes m_i}
\end{align}
of $\U$-crystals, 
which will be viewed as a counterpart of~\eqref{eq:emb-sst} in this paper 
(see~\eqref{eq:LT} and Lemma~\ref{lem:Phi_LT}). 
In the case that $\U$ is of untwisted affine type $A$, 
it follows that 
$\mathcal{B}(\varpi_i)$
is parametrized by 
$\CST (\varpi_i) \times \BZ$, 
and the affine Weyl group acts transitively on 
$\CST (\varpi_i) \times \BZ$ 
(see Proposition~\ref{prop:aff-varpi_i}). 
Consequently, 
$\bigotimes_{i \in I} \mathcal{B}(\varpi_i)^{\otimes m_i}$ 
is parametrized by 
$\CST (\lambda) \times \BZ^N$, 
where $N := \sum_{i \in I} m_i$, 
and each element in 
$\CST (\lambda) \times \BZ^N$
is labeled by a tuple of $N$ cosets in the affine Weyl group. 
Let $\mathbb{T} \in \CST (\lambda) \times \BZ^N$, 
and assume that $\mathbb{T}$ corresponds to the tuple 
$(\bar{x}_1 , \bar{x}_2 , \ldots , \bar{x}_N)$
of cosets in the affine Weyl group. 
Then 
$\mathbb{T}$ 
is in the image of the strict embedding~\eqref{eq:intro-LT} 
if and only if there exist coset representatives 
$x_{\nu} \in \bar{x}_{\nu}$, $1 \leq \nu \leq N$, 
such that
$x_1 \succeq x_2 \succeq \cdots \succeq x_N$ 
in the semi-infinite Bruhat order $\succeq$ on the affine Weyl group
(see Theorem~\ref{thm:SMT}). 
Such decreasing sequences are explicitly described in terms of tableaux 
(see Definition~\ref{def:SiYT}~\ref{Defi4.2_1} and Theorem~\ref{thm:tab-cri}). 
We are thus led to the definition of semi-infinite Young tableaux 
(see Definition~\ref{def:SiYT}~\ref{Defi4.2_2}).

Let us give an example of a semi-infinite Young tableau
(by using the notation in \S~\ref{subsection:SiYT}). 
Let $n = 7$ and 
\begin{align}
\lambda 
= 
\varpi_1 + 3 \varpi_2 + 2\varpi_4
=
\begin{ytableau}
\ & \ & \ & \ & \ & \ \\
\none & \ & \ & \ & \ & \ \\
\none & \none & \none & \none & \ & \ \\
\none & \none & \none & \none & \ & \ 
\end{ytableau} \, .
\end{align}
We claim that 
\begin{align}
\mathbb{T}
=
\left(
\begin{ytableau}
5 & 2 & 3 & 1 & 1 & 1 \\
\none & 3 & 4 & 5 & 3 & 4 \\
\none & \none & \none & \none & 4 & 5 \\
\none & \none & \none & \none & 6 & 7 
\end{ytableau}
\, , \ 
(-3;5,4,-1;8,7)
\right)
\in 
\CST (\lambda) \times \BZ^6
\end{align}
is a semi-infinite Young tableau of shape $\lambda$. 
Indeed, 
$\mathbb{T}$
is a semi-infinite Young tableau if (and only if)
its rectangle components
\begin{align}
\left(
\begin{ytableau}
5
\end{ytableau}
\, , \ 
-3
\right), 
\hspace{5mm}
\left(
\begin{ytableau}
2 & 3 & 1 \\ 
3 & 4 & 5
\end{ytableau}
\, , \ 
(5,4,-1) 
\right), 
\hspace{5mm}
\left(
\begin{ytableau}
1 & 1 \\ 
3 & 4 \\
4 & 5 \\
6 & 7
\end{ytableau}
\, , \ 
(8,7)
\right)
\end{align}
are semi-infinite Young tableaux. 
Also, a tableau
{\postdisplaypenalty-1000
\begin{align}
(\mathsf{T}_1 \mathsf{T}_2 \cdots \mathsf{T}_m , (c_1 , c_2 , \ldots , c_m))
\in \CST (m\varpi_i) \times \BZ^m
\end{align}}%
of rectangle shape is a semi-infinite Young tableau 
if (and only if)
\begin{enumerate}[label=(\roman*)]
\item
$c_1 \geq c_2 \geq \cdots \geq c_m$, 
and
\item
$\bigl( 
\mathsf{T}_{\nu}(u + c_{\nu} - c_{\nu+1}) \geq \mathsf{T}_{\nu+1}(u)$
if 
$1 \leq u \leq i - c_{\nu} + c_{\nu+1}
\bigr)$
for every
$1 \leq \nu < m$, 
\end{enumerate}
where 
$\mathsf{T}_{\nu}(s)$
denotes the $s$-th entry
(from top) of the $\nu$-th column
$\mathsf{T}_{\nu}$.

This paper is organized as follows. 
In \S~\ref{Section:Weyl}, 
we set up notation and terminology on 
untwisted affine root data and crystals. 
Also, we have compiled some basic facts on 
extremal weight modules over quantized universal enveloping algebras of untwisted affine types, 
the semi-infinite Bruhat order on affine Weyl groups, 
and semi-infinite Lakshmibai--Seshadri paths. 
In \S~\ref{Section:SMT}, 
we introduce a strict embedding 
$\Phi_{\lambda \mid q=0}^{\LT }$ 
of crystals. 
Then we state and prove a characterization of the image of 
$\Phi_{\lambda \mid q=0}^{\LT }$, 
which can be thought of as standard monomial theory for 
semi-infinite Lakshmibai--Seshadri paths (see Theorem~\ref{thm:SMT}). 
In \S~\ref{Section:SiYT}, 
we will restrict our attention to the case of 
untwisted affine type $A$. 
We introduce semi-infinite Young tableaux. 
We prove that the set 
$\mathbb{Y}^{\si}(\lambda)$ 
of semi-infinite Young tableaux of shape $\lambda$ 
equals the image of 
$\Phi_{\lambda \mid q=0}^{\LT }$
(see Theorem~\ref{thm:SiYT}), 
by showing a tableau criterion for the semi-infinite Bruhat order 
(see Theorem~\ref{thm:tab-cri}). 
Consequently, this proves that 
$\mathbb{Y}^{\si}(\lambda)$ 
is isomorphic, as a $\U$-crystal, to the crystal basis 
$\mathcal{B}(\lambda)$ 
(see Corollary~\ref{cor:isom}). 
We give an explicit description of the crystal structure on 
$\mathbb{Y}^{\si}(\lambda)$ 
(see Proposition~\ref{prop:crystal-str-SiYT}). 

\section{Preliminaries} \label{Section:Weyl}

\subsection{Untwisted affine root data}

Let 
$\Fg_{\af}$ 
be an untwisted affine Lie algebra over $\BC$ 
with a Cartan subalgebra $\Fh_{\af}$. 
Let 
$\{ \alpha_i \}_{i \in I_{\af}} \subset \Fh_{\af}^* = \Hom_{\BC}(\Fh_{\af},\BC)$ 
and 
$\{ h_i \}_{i \in I_{\af}} \subset \Fh_{\af}$
be the sets of simple roots and simple coroots, respectively. 
Here $I_{\af}$ denotes the vertex set of the (affine) Dynkin diagram of $\Fg_{\af}$. 
Let 
$\langle \cdot , \cdot \rangle : \Fh_{\af} \times \Fh_{\af}^* \rightarrow \BC$ 
be the canonical pairing. 
We take and fix an integral weight lattice 
$P_{\af} \subset \Fh_{\af}^*$ 
satisfying the conditions that 
$\alpha_i \in P_{\af}$ 
and 
$h_i \in \Hom_{\BZ}(P_{\af},\BZ)$ 
for all 
$i \in I_{\af}$, and for each 
$i \in I_{\af}$
there exists 
$\Lambda_i \in P_{\af}$ 
such that 
$\langle h_j , \Lambda_i \rangle = \delta_{ij}$ 
for all 
$j \in I_{\af}$. 
Let 
$\delta = \sum_{i \in I_{\af}} a_i \alpha_i \in \Fh_{\af}^*$ 
and 
$c = \sum_{i \in I_{\af}} a_i^{\vee} h_i \in \Fh_{\af}$ 
be the null root and the canonical central element, respectively.
We take and fix $0 \in I_{\af}$ such that 
$a_0 = a_0^{\vee} = 1$. 
Set 
$I = I_{\af} \setminus \{ 0 \}$; 
note that the subset $I$ of $I_{\af}$ 
corresponds to the vertex set of the Dynkin diagram of a 
complex finite-dimensional simple Lie subalgebra $\Fg$ of $\Fg_{\af}$. 
For each 
$i \in I_{\af}$, 
define
$\varpi_i = \Lambda_i - \langle c , \Lambda_i \rangle \Lambda_0$
and call it the 
$i$-th level-zero fundamental weight; 
note that 
$\varpi_0 = 0$, 
$\langle c , \varpi_i \rangle = 0$
for all $i \in I_{\af}$, and 
$\langle h_i , \varpi_j \rangle = \delta_{ij}$ 
for all $i,j \in I$.

Let 
$W_{\af} = \langle r_i \mid i \in I_{\af} \rangle$ 
be the (affine) Weyl group of $\Fg_{\af}$, 
where $r_i$ denotes the simple reflection with respect to $\alpha_i$.
The subgroup 
$W = \langle r_i \mid i \in I \rangle \subset W_{\af}$ 
is the (finite) Weyl group of $\Fg$. 
Let $\ell : W_{\af} \rightarrow \BZ_{\geq 0}$ be the length function. 
Let $e \in W_{\af}$ be the unit element. 
Let 
$(\cdot , \cdot)$
be a $W_{\af}$-invariant non-degenerate symmetric bilinear form on $\Fh_{\af}^*$ such that
$(\delta,\lambda) = \langle c,\lambda \rangle$ 
for all $\lambda \in \Fh_{\af}^*$. 
Set $\alpha_i^{\vee} = \dfrac{2 \alpha_i}{(\alpha_i , \alpha_i)} \in \Fh_{\af}^*$ for $i \in I_{\af}$.
The action of $W_{\af}$ on $\Fh_{\af}^*$ is given by
$r_i (\lambda) 
= 
\lambda - (\alpha_i^{\vee} , \lambda) \alpha_i
=
\lambda - \langle h_i , \lambda \rangle \alpha_i$ 
for $i \in I_{\af}$ and $\lambda \in \Fh_{\af}^*$.
Set
\begin{align}
Q
= 
\bigoplus_{i \in I} \BZ \alpha_i, &&
Q^{\vee} 
= 
\bigoplus_{i \in I} \BZ \alpha_i^{\vee}, &&
P^+ 
= 
\sum_{i \in I} \BZ_{\geq 0} \varpi_i.
\end{align}
We know from~\cite[\S~6.5]{Kac} that 
$Q^{\vee} \subset Q$. 
For $\xi \in Q^{\vee}$, 
we denote by 
$t_{\xi} \in W_{\af}$ 
the translation by $\xi$ (see~\cite[\S6.5]{Kac}).
We know from~\cite[Proposition~6.5]{Kac} 
that 
$\{ t_{\xi} \mid \xi \in Q^{\vee} \}$ 
forms an abelian normal subgroup of $W_{\af}$, 
for which 
$t_{\xi} t_{\zeta} = t_{\xi + \zeta}$, 
$\xi , \zeta \in Q^{\vee}$, 
and 
$W_{\af} = W \ltimes \{ t_{\xi} \mid \xi \in Q^{\vee} \}$. 
For $w \in W$ and $\xi \in Q^{\vee}$, we have 
\begin{align} \label{eq:lv0-act}
wt_{\xi}\lambda = w\lambda - ( \xi,\lambda ) \delta \ \ 
\text{if} \ 
\lambda \in \Fh_{\af}^* \ 
\text{satisfies} \ 
\langle c,\lambda \rangle = 0.
\end{align}

Let $\Delta$ be the root system of $\Fg$. 
Set $\Delta^+ = \Delta \cap \sum_{i \in I} \BZ_{\geq 0} \alpha_i$. 
For a subset $J \subset I$, set
\begin{align}
Q_J = \bigoplus_{j \in J} \BZ \alpha_j, && 
Q_J^{\vee} = \bigoplus_{j \in J} \BZ \alpha_j^{\vee}, && 
\Delta_J = \Delta \cap Q_J, && 
\Delta_J^+ = \Delta^+ \cap Q_J. 
\end{align}
Denote by $\Delta_{\af}$ the set of real roots of $\Fg_{\af}$, 
and by $\Delta_{\af}^+$ the set of positive real roots of $\Fg_{\af}$; 
we know from~\cite[Proposition~6.3]{Kac} that
\begin{align}
\Delta_{\af} 
= 
\{ \alpha + n\delta \mid \alpha \in \Delta , n \in \BZ \} , &&
\Delta_{\af}^+
= 
\Delta^+ 
\sqcup
\{ \alpha + n\delta \mid \alpha \in \Delta , n \in \BZ_{>0} \} .
\end{align}
For $\beta \in \Delta_{\af}$, 
let 
$\beta^{\vee} = \dfrac{2\beta}{(\beta,\beta)} \in \Fh_{\af}^*$, 
and let 
$r_{\beta} \in W_{\af}$ 
be the reflection with respect to $\beta$; 
if $\beta = \alpha + n\delta$, $\alpha \in \Delta$ and $n \in \BZ$, then
$r_{\beta} = r_{\alpha} t_{n\alpha^{\vee}}$.

Let $d$ be the smallest positive integer such that
$(\alpha_i , \alpha_i) / 2 \in (1/d)\BZ$ 
for all $i \in I_{\af}$. 
Let $q$ be an indeterminate, and set
$q_s = q^{1/d}$. 
Let 
$\U$
be the quantized universal enveloping algebra over $\BQ(q_s)$
associated with $\Fg_{\af}$. 
Let 
$\U'$
be the $\BQ(q_s)$-subalgebra of $\U$ 
corresponding to the derived subalgebra
$[\Fg_{\af},\Fg_{\af}]$ of $\Fg_{\af}$
(see for instance~\cite[\S~2.2]{BN04}). 

\subsection{Crystals} \label{Subsection:Crystals}

In this subsection, 
we set up notation and terminology on crystals. 
For a fuller treatment, we refer the reader to~\cite{AK,K94,K95,K02}. 

A set
$\mathcal{B}$
together with the maps
$\wt : \mathcal{B} \rightarrow P_{\af}$
(\resp $\wt : \mathcal{B} \rightarrow P_{\af} / (P_{\af} \cap \BC \delta)$), 
$e_i , f_i : \mathcal{B} \rightarrow \mathcal{B} \sqcup \{ \bm{0} \}$, $i \in I_{\af}$, 
and 
$\varepsilon_i , \varphi_i : \mathcal{B} \rightarrow \BZ \sqcup \{ -\infty \}$
is called a $\U$-crystal (\resp $\U'$-crystal) if the following conditions are satisfied:
\begin{enumerate}[label=(C\arabic*)]
\item
$\varphi_i(b) = \varepsilon_i(b) + \langle h_i , \wt (b) \rangle$ 
for all 
$i \in I_{\af}$, 
\item
$\wt (e_i b) = \wt (b) + \alpha_i$ 
if 
$e_i b \in \mathcal{B}$, 
\item
$\wt (f_i b) = \wt (b) - \alpha_i$ 
if 
$f_i b \in \mathcal{B}$, 
\item
$\varepsilon_i (e_i b) = \varepsilon_i(b) - 1$
and 
$\varphi_i (e_i b) = \varphi_i(b) + 1$
if 
$e_i b \in \mathcal{B}$, 
\item
$\varepsilon_i (f_i b) = \varepsilon_i(b) + 1$
and 
$\varphi_i (f_i b) = \varphi_i(b) - 1$
if 
$f_i b \in \mathcal{B}$, 
\item
$f_i b = b'$
if and only if 
$b = e_i b'$
for 
$b,b' \in \mathcal{B}$ 
and 
$i \in I_{\af}$, 
\item
if 
$\varphi_i(b) = -\infty$, 
then
$e_i b = f_i b = \bm{0}$.
\end{enumerate}
The maps 
$e_i , f_i : \mathcal{B} \rightarrow \mathcal{B} \sqcup \{ \bm{0} \}$, $i \in I_{\af}$, 
are called the Kashiwara operators. 
For a subset $\mathcal{B}'$ of a crystal $\mathcal{B}$, 
we say that $\mathcal{B}'$ is stable under the Kashiwara operators if
$e_i \mathcal{B}' , f_i \mathcal{B}' \subset \mathcal{B}' \sqcup \{ \bm{0} \}$
for all 
$i \in I_{\af}$.

Let 
$\mathcal{B}_1$
and 
$\mathcal{B}_2$
be $\U$-crystals or $\U'$-crystals.
A morphism 
$\Psi : \mathcal{B}_1 \rightarrow \mathcal{B}_2$
is, by definition, a map
$\mathcal{B}_1 \sqcup \{ \bm{0} \} \rightarrow \mathcal{B}_2 \sqcup \{ \bm{0} \}$
such that 
\begin{enumerate}[label=(M\arabic*)]
\item
$\Psi (\bm{0}) = \bm{0}$, 
\item
if 
$b \in \mathcal{B}_1$ 
and 
$\Psi(b) \in \mathcal{B}_2$, 
then 
$\wt (\Psi(b)) = \wt (b)$, 
$\varepsilon_i (\Psi(b)) = \varepsilon_i (b)$, and
$\varphi_i(\Psi(b)) = \varphi_i(b)$ 
for all 
$i \in I_{\af}$, 
\item
if
$b , b' \in \mathcal{B}_1$, 
$\Psi(b) , \Psi(b') \in \mathcal{B}_2$ and 
$f_i b = b'$, 
then
$f_i \Psi(b) = \Psi(b')$
for all 
$i \in I_{\af}$. 
\end{enumerate}
A morphism
$\Psi : \mathcal{B}_1 \rightarrow \mathcal{B}_2$
is called strict if 
$\Psi(f_i b) = f_i \Psi (b)$
and 
$\Psi(e_i b) = e_i \Psi (b)$
for all 
$b \in \mathcal{B}_1$ and $i \in I_{\af}$.
A morphism
$\Psi : \mathcal{B}_1 \rightarrow \mathcal{B}_2$
is called a strict embedding if it is a strict morphism and the associated map 
$\mathcal{B}_1 \sqcup \{ \bm{0} \} \rightarrow \mathcal{B}_2 \sqcup \{ \bm{0} \}$
is injective. A morphism
$\Psi : \mathcal{B}_1 \rightarrow \mathcal{B}_2$
is called an isomorphism if the associated map 
$\mathcal{B}_1 \sqcup \{ \bm{0} \} \rightarrow \mathcal{B}_2 \sqcup \{ \bm{0} \}$
is bijective.

The tensor product
$\mathcal{B}_1 \otimes \mathcal{B}_2$ 
of crystals
$\mathcal{B}_1$
and 
$\mathcal{B}_2$ 
is defined to be the set 
$\mathcal{B}_1 \times \mathcal{B}_2$ 
whose crystal structure is defined as follows: 
\begin{enumerate}[label=(T\arabic*)]
\item
$\wt (b_1 \otimes b_2) = \wt (b_1) + \wt (b_2)$, 
\item
$\varepsilon_i (b_1 \otimes b_2) 
= 
\max \{ \varepsilon_i(b_1) , \varepsilon_i(b_2) - \langle h_i , \wt (b_1) \rangle \}$, 
\item
$\varphi_i (b_1 \otimes b_2) 
= 
\max \{ \varphi_i(b_2) , \varphi_i(b_1) + \langle h_i , \wt (b_2) \rangle \}$, 
\item
$e_i (b_1 \otimes b_2) 
=
\begin{cases}
(e_i b_1) \otimes b_2 
&\text{if} \ \varphi_i(b_1) \geq \varepsilon_i(b_2), \\
b_1 \otimes (e_i b_2) 
&\text{if} \ \varphi_i(b_1) < \varepsilon_i(b_2),
\end{cases}$
\item
$f_i (b_1 \otimes b_2) 
=
\begin{cases}
(f_i b_1) \otimes b_2 
&\text{if} \ \varphi_i(b_1) > \varepsilon_i(b_2), \\
b_1 \otimes (f_i b_2) 
&\text{if} \ \varphi_i(b_1) \leq \varepsilon_i(b_2).
\end{cases}$
\end{enumerate}
Here, we write 
$b_1 \otimes b_2$ 
for 
$(b_1 , b_2) \in \mathcal{B}_1 \times \mathcal{B}_2$, 
and we understand that
$b_1 \otimes \bm{0} = \bm{0} \otimes b_2 = \bm{0}$. 

Let 
$\mathcal{B}$ 
be a regular $\U$-crystal (\resp regular $\U'$-crystal) in the sense of~\cite[\S~2.2]{K02}. 
By~\cite[\S~7]{K94}, 
we have a $W_{\af}$-action 
$S : W_{\af} \rightarrow \Bij(\mathcal{B})$, $x \mapsto S_x$, 
on (the underlying set) 
$\mathcal{B}$
given by
\begin{align} \label{eq:W-action}
S_{r_i} b
=
\begin{cases}
f_i^{\langle h_i , \wt (b) \rangle} b & \text{if}\ \langle h_i , \wt (b) \rangle \geq 0, \\%[1mm]
e_i^{-\langle h_i , \wt (b) \rangle} b & \text{if}\ \langle h_i , \wt (b) \rangle \leq 0, 
\end{cases}
\end{align}
for each 
$b \in \mathcal{B}$ 
and 
$i \in I_{\af}$. 
Note that 
$\wt (S_x b) = x \wt (b)$ 
holds for all
$x \in W_{\af}$ and $b \in \mathcal{B}$. 
An element 
$b \in \mathcal{B}$ 
of weight 
$\lambda \in P_{\af}$
(\resp $\lambda \in P_{\af} / (P_{\af} \cap \BC \delta)$)
is called an extremal element if we can find elements
$b_x \in \mathcal{B}$, $x \in W_{\af}$, 
such that 
\begin{enumerate}[label=(E\arabic*)]
\item
$b_e = b$, 
\item
if 
$\langle h_i , x\lambda \rangle \geq 0$, 
then
$e_i b_x = \bm{0}$
and 
$f_i^{\langle h_i , x\lambda \rangle} b_x = b_{r_i x}$, 
\item
if 
$\langle h_i , x\lambda \rangle \leq 0$, 
then
$f_i b_x = \bm{0}$
and 
$e_i^{-\langle h_i , x\lambda \rangle} b_x = b_{r_i x}$. 
\end{enumerate}
Then 
$b_x = S_x b$
holds for all 
$x \in W_{\af}$.

\subsection{Extremal weight modules and their crystal bases} \label{Subsection:extremal}

In this subsection, 
following~\cite{BN04,K94,K02}, 
we review some of the standard facts on 
extremal weight modules and their crystal bases.

For 
$\lambda \in P^+$, 
let 
$V(\lambda)$ 
be the extremal weight $\U$-module generated by an extremal weight vector 
$u_{\lambda}$ 
of extremal weight 
$\lambda$, 
and let 
$(\mathcal{L}(\lambda) , \mathcal{B}(\lambda))$ 
be the crystal basis of 
$V(\lambda)$
(\cite[Proposition~8.2.2]{K94}; see also~\cite[\S~3.2]{K02}). 
Note that 
$\mathcal{B}(\lambda)$ 
is a regular $\U$-crystal in the sense of~\cite[\S~2.2]{K02} (see \S~\ref{Subsection:Crystals}). 
Let 
$\mathcal{B}_0(\lambda)$ 
be the connected component of the crystal graph of 
$\mathcal{B}(\lambda)$
containing 
$u_{\lambda} \mod q_s \mathcal{L}(\lambda)$.

Let $z_i$, $i \in I$, be the $\U'$-linear automorphism of $V(\varpi_i)$ of weight $\delta$ introduced in~\cite[\S~5.2]{K02}; $z_i$ sends a (unique) global basis element of weight $\varpi_i$ to a (unique) global basis element of weight $\varpi_i+\delta$. 
Then $z_i$ induces a
$\BQ$-linear automorphism of 
$\mathcal{L}(\varpi_i) / q_s \mathcal{L}(\varpi_i)$
and an automorphism of 
$\mathcal{B}(\varpi_i)$ 
as a $\U'$-crystal; 
by abuse of notation, 
we use the same letter 
$z_i$ 
for the automorphism of 
$\mathcal{B}(\varpi_i)$. 
The $\U'$-module 
$W(\varpi_i) = V(\varpi_i) / (z_i - 1)V(\varpi_i)$
is called a level-zero fundamental representation. 
We know from~\cite[Theorem~5.17]{K02} that 
$W(\varpi_i)$ is a finite-dimensional irreducible $\U'$-module 
and has a (simple) crystal basis. 

For 
$\lambda = \sum_{i \in I} m_i \varpi_i \in P^+$, 
with 
$m_i \in \BZ_{\geq 0}$, $i \in I$, 
let 
$\Tilde{V}(\lambda) = \bigotimes_{i \in I} V(\varpi_i)^{\otimes m_i}$ 
and 
$\Tilde{u}_{\lambda} = \bigotimes_{i \in I} u_{\varpi_i}^{\otimes m_i} \in \Tilde{V}(\lambda)$. 
For each $i \in I$ and $1 \leq \nu \leq m_i$, let $z_{i,\nu}$ be the $\U'$-linear automorphism of $\Tilde{V}(\lambda)$ obtained by the action of $z_i$ on the $\nu$-th factor $V(\varpi_i)$ of $V(\varpi_i)^{\otimes m_i}$ in $\Tilde{V}(\lambda)$. The $\U$-submodule 
\begin{align}
\Breve{V}(\lambda) = \U \left[z_{i,\nu}, z_{i,\nu}^{-1} \ \vline \ i \in I , \, 1 \leq \nu \leq m_i \right] \Tilde{u}_{\lambda} \subset \Tilde{V}(\lambda)
\end{align}
has a crystal basis 
$(\Breve{\mathcal{L}}(\lambda) , \Breve{\mathcal{B}}(\lambda) = \bigotimes_{i \in I} \mathcal{B}(\varpi_i)^{\otimes m_i})$ 
such that 
$\Breve{\mathcal{L}}(\lambda) \subset \bigotimes_{i \in I} \mathcal{L}(\varpi_i)^{\otimes m_i}$
(\cite[Theorem~8.5]{K02}). 
Let
$\Breve{\mathcal{B}}_0(\lambda)$
be the connected component of the crystal graph of 
$\Breve{\mathcal{B}}(\lambda)$
containing 
$\Tilde{u}_{\lambda} \mod q_s \Breve{\mathcal{L}}(\lambda)$.
Since $\Breve{V}(\lambda)$ contains an extremal weight vector 
$\Tilde{u}_{\lambda}$ of weight $\lambda$, 
we have a $\U$-linear homomorphism
$\Phi_{\lambda} : V(\lambda) \rightarrow \Breve{V}(\lambda)$
sending $u_{\lambda}$ to $\Tilde{u}_{\lambda}$. 
We know from~\cite[\S~4.2]{BN04} that the map 
$\Phi_{\lambda}$ 
is injective, 
commutes with the Kashiwara operators 
$e_i , f_i$, $i \in I_{\af}$, 
and induces an injective $\BQ$-linear map
$\Phi_{\lambda \mid q=0} : \mathcal{L}(\lambda) / q_s \mathcal{L}(\lambda) \rightarrow \Breve{\mathcal{L}}(\lambda) / q_s \Breve{\mathcal{L}}(\lambda)$; 
note that 
$\Phi_{\lambda \mid q=0}(\mathcal{B}(\lambda)) \not\subset \Breve{\mathcal{B}}(\lambda)$, 
in general (see Theorem~\ref{thm:BN}~\ref{Theo2.1_2}).

For 
$\lambda = \sum_{i \in I} m_i \varpi_i \in P^+$, 
with 
$m_i \in \BZ_{\geq 0}$, $i \in I$, 
set 
\begin{align} \label{eq:Par}
\Par (\lambda) = \left\{ \bmrho{\rho} = \left( \rho^{(i)} \right)_{i \in I} \ \vline \ 
\rho^{(i)} \ \text{is a partition of length less than} \ m_i \ \text{for each} \ i \in I \right\} ;
\end{align}
\looseness-1
we understand that a partition of length less than $1$ is an empty partition $\emptyset$. 
For a partition $\rho = (\rho_1 \geq \rho_2 \geq \cdots \geq \rho_l > 0)$, 
set 
$|\rho| = \sum_{\nu = 1}^l \rho_{\nu}$.
For 
$\bmrho{\rho} = \left( \rho^{(i)} \right)_{i \in I} \in \Par (\lambda)$, 
set 
\begin{align}
\wt (\bmrho{\rho}) = -\sum_{i \in I} \left| \rho^{(i)} \right| \delta .
\end{align}

Let $\bmrho{\rho} = \left( \rho^{(i)} \right)_{i \in I} \in \Par (\lambda)$. 
Let
$S_{\bmrho{\rho}}^-$
be the (PBW-type) basis element 
of weight $\wt (\bmrho{\rho})$ 
of the negative imaginary part of $\U$ 
constructed in~\cite[the element $S_{\mathbf{c}_0}^-$ in \S~3.1; see also Remark~4.1]{BN04}.
Define
\begin{align} \label{eq:Schur}
s_{\bmrho{\rho}}(z^{-1})
=
\prod_{i \in I} s_{\rho^{(i)}} ( z_{i,1}^{-1} , z_{i,2}^{-1} , \ldots , z_{i,m_i}^{-1} ),
\end{align}
where the right-hand side is a product of Schur polynomials in the variables 
$z_{i,\nu}^{-1}$, $i \in I$, $1 \leq \nu \leq m_i$.

\begin{theo}[{\cite[\S~4.2]{BN04}}; see also {\cite[\S~13]{K02}}] \label{thm:BN}
Let $\lambda \in P^+$. 
\begin{enumerate}[label=(\arabic*)]
\item\label{Theo2.1_1}
We have
\begin{equation}
\begin{aligned}
%\begin{split}
\mathcal{B}(\lambda)
=
\Bigl\{ 
&g_1 g_2 \cdots g_l S_{\bmrho{\rho}}^- u_{\lambda} 
\mod q_s \mathcal{L}(\lambda)
 \\%[1mm]
&
\Bigl|\ g_k \in \{ e_i , f_i \mid i \in I_{\af} \} , \ 
1 \leq k \leq l , \ 
l \in \BZ_{\geq 0}, \ 
\bmrho{\rho} \in \Par (\lambda) 
\Bigr\} 
\, \setminus \, 
\left\{ 0 \right\} .
%\end{split}
\end{aligned}
\end{equation}
\item\label{Theo2.1_2}
We have 
\begin{equation}
\label{eq:Phi_q=0}
\begin{aligned} 
%\begin{split}
\Phi_{\lambda \mid q=0} \left( g_1 g_2 \cdots g_l S_{\bmrho{\rho}}^- u_{\lambda} 
\mod q_s \mathcal{L}(\lambda) \right) 
\, &= \, 
g_1 g_2 \cdots g_l s_{\bmrho{\rho}}(z^{-1}) \Tilde{u}_{\lambda} 
\mod q_s \Breve{\mathcal{L}}(\lambda) \\%[1mm]
&= 
s_{\bmrho{\rho}}(z^{-1}) g_1 g_2 \cdots g_l \Tilde{u}_{\lambda} 
\mod q_s \Breve{\mathcal{L}}(\lambda)
%\end{split}
\end{aligned}
\end{equation}
for 
$g_1 g_2 \cdots g_l S_{\bmrho{\rho}}^- u_{\lambda} 
\mod q_s \mathcal{L}(\lambda)
\in \mathcal{B}(\lambda)$. 
In particular, 
\begin{align}
\Phi_{\lambda \mid q=0} (\mathcal{B}(\lambda))
=
\left\{ 
s_{\bmrho{\rho}}(z^{-1})b 
\mod q_s \Breve{\mathcal{L}}(\lambda)
\ \vline \ 
\bmrho{\rho} \in \Par (\lambda) , \ 
b \in \Breve{\mathcal{B}}_0(\lambda) 
\right\} ,
\end{align}
and the map 
$\Phi_{\lambda \mid q=0}$
induces an isomorphism of $\U$-crystals from
$\mathcal{B}_0(\lambda)$
to $\Breve{\mathcal{B}}_0(\lambda)$.
\item\label{Theo2.1_3}
Let 
$\mathcal{B}_{\bmrho{\rho}}(\lambda)$
be the connected component of $\mathcal{B}(\lambda)$ containing
$S_{\bmrho{\rho}}^- u_{\lambda} \mod q_s \mathcal{L}(\lambda)$. 
Then we have 
$\mathcal{B}(\lambda) = \bigsqcup_{\bmrho{\rho} \in \Par (\lambda)} \mathcal{B}_{\bmrho{\rho}}(\lambda)$.
Moreover, 
for each 
$\bmrho{\rho} \in \Par (\lambda)$, 
there exists an isomorphism of $\U'$-crystals from
$\mathcal{B}_0(\lambda)$
to
$\mathcal{B}_{\bmrho{\rho}}(\lambda)$
sending
$u_{\lambda} \mod q_s \mathcal{L}(\lambda)$
to
$S_{\bmrho{\rho}}^- u_{\lambda} \mod q_s \mathcal{L}(\lambda)$.
\end{enumerate}
\end{theo}

\subsection{Semi-infinite Bruhat order on affine Weyl groups} \label{Subsection:SiBO}

In this subsection, 
we recall some basic facts on 
the semi-infinite Bruhat order on affine Weyl groups 
(see~\cite{INS16,LS,P} for more details).

We take and fix $J \subset I$. 
Let $W_J = \langle r_j \mid j \in J \rangle$, 
and let 
$W^J$
be the set of minimal coset representatives for $W/W_J$ 
(see~\cite[Corollary~2.4.5~(i)]{BB}). 
For 
$w \in W$, 
we denote by 
$\lfloor w \rfloor \in W^J$ 
the minimal coset representative for the coset 
$wW_J \in W/W_J$.
Define
{\postdisplaypenalty1000000
\begin{align}
(\Delta_J)_{\af} 
&=
\{ \alpha + n\delta \mid \alpha \in \Delta_J , n \in \BZ \} \subset \Delta_{\af}, \\%[1mm]
(\Delta_J)_{\af}^+ 
&=
(\Delta_J)_{\af} \cap \Delta_{\af}^+ 
=
\Delta_J^+ 
\sqcup
\{ \alpha + n\delta \mid \alpha \in \Delta_J , n \in \BZ_{>0} \}, \\%[1mm]
(W_J)_{\af}
&=
W_J \ltimes \{ t_{\xi} \mid \xi \in Q_J^{\vee} \}
=
\langle r_{\beta} \mid \beta \in (\Delta_J)_{\af}^+ \rangle , \label{eq:W_J_af} \\%[1mm]
(W^J)_{\af} 
&=
\{ x \in W_{\af} \mid x\beta \in \Delta_{\af}^+ \ \text{for all} \ \beta \in (\Delta_J)_{\af}^+ \} ;
\end{align}}%
note that 
$(W_{\emptyset})_{\af} = \{ e \}$ 
and 
$(W^{\emptyset})_{\af} = W_{\af}$.

\begin{lemm}[{\cite{P}; see also~\cite[Lemma~10.5]{LS}}] \label{lem:Waf-decomp}
For each $x \in W_{\af}$, 
there exist a unique $x_1 \in (W^J)_{\af}$
and a unique $x_2 \in (W_J)_{\af}$ such that
$x = x_1 x_2$. 
In particular, 
$(W^J)_{\af}$ 
is a complete system of coset representatives for 
$W_{\af} / (W_J)_{\af}$. 
\end{lemm}

Define a map 
$\Pi^J : W_{\af} \rightarrow (W^J)_{\af}$ 
by $\Pi^J(x) = x_1$ 
if $x = x_1 x_2$ with 
$x_1 \in (W^J)_{\af}$
and 
$x_2 \in (W_J)_{\af}$. 

\begin{lemm}[{\cite{P}; see also~\cite[Proposition~10.8]{LS}}] \label{lem:Pi}\ \\[-1.2em]

\begin{enumerate}[label=(\arabic*)]
\item\label{Lemm2.3_1}
$\Pi^J(w) = \lfloor w \rfloor$ 
for $w \in W$.
\item \label{Lemm2.3_2}
$\Pi^J(x t_{\xi}) = \Pi^J(x) \Pi^J(t_{\xi})$
for $x \in W_{\af}$ and $\xi \in Q^{\vee}$.
\end{enumerate}
\end{lemm}

For simplicity of notation, we let $T_{\xi} = T_{\xi}^J$ stand for 
$\Pi^J(t_{\xi}) \in (W^J)_{\af}$ for $\xi \in Q^{\vee}$. 
The next lemma follows immediately from~\eqref{eq:W_J_af} and Lemmas~\ref{lem:Waf-decomp}--\ref{lem:Pi}.

\begin{lemm} \label{lem:T_xi}\ \\[-1.2em]
\begin{enumerate}[label=(\arabic*)]
\item\label{Lemm2.4_1}
$(W^J)_{\af} = \{ wT_{\xi} \mid w \in W^J , \ \xi \in Q^{\vee} \}$.
\item\label{Lemm2.4_2}
Let 
$\xi , \zeta \in Q^{\vee}$. 
If 
$\xi \equiv \zeta \mod Q_J^{\vee}$, 
then 
$T_{\xi}^J = T_{\zeta}^J$. 
\end{enumerate}
\end{lemm}

Set $\rho_J = (1/2)\sum_{\alpha \in \Delta_J^+} \alpha$; 
we abbreviate $\rho_J$ to $\rho$ if $J = I$. 
For $x = wt_{\xi} \in W_{\af}$ with $w \in W$ and $\xi \in Q^{\vee}$, define
\begin{align}
\ell^{\si}(x) = \ell(w) + 2 ( \xi,\rho ).
\end{align}

Define the (parabolic) semi-infinite Bruhat graph $\SiB ^J$ 
to be the $\Delta_{\af}^+$-colored directed graph 
with vertex set $(W^J)_{\af}$ and edges of the form 
$x \xrightarrow{\ \beta \ } r_{\beta} x$ 
for 
$x \in (W^J)_{\af}$ and $\beta \in \Delta_{\af}^+$, 
where
$r_{\beta} x \in (W^J)_{\af}$ 
and 
$\ell^{\si}(r_{\beta}x) = \ell^{\si}(x) + 1$.

The semi-infinite Bruhat order is a partial order $\preceq$ on $(W^J)_{\af}$ 
defined as follows: 
for $x , y \in (W^J)_{\af}$, we write $x \preceq y$ 
if there exists a directed path from $x$ to $y$ in $\SiB ^J$.

\begin{prop}[{\cite[Proposition~A.1.2]{INS16}}] \label{prop:phi_J}
Let $w \in W^J$, $\xi \in Q^{\vee}$ and $\beta \in \Delta_{\af}^+$. 
Write 
$\beta = w\gamma + \chi\delta$ 
with 
$\gamma \in \Delta$ 
and 
$\chi \in \BZ_{\geq 0}$.
Then 
$r_{\beta}wT_{\xi} \in (W^J)_{\af}$
and there exists an edge 
$wT_{\xi} \xrightarrow{\ \beta \ } r_{\beta}wT_{\xi}$ in $\SiB ^J$ 
if and only if $\gamma \in \Delta^+ \setminus \Delta_J^+$ 
and one of the following conditions holds:
\begin{enumerate}[label=(\Alph*)]\setcounter{enumi}{1}
\item\label{Prop2.5_B} %[(B)]
$\chi = 0$
and 
$\ell(w r_{\gamma}) = \ell(w) + 1$; 
in this case, we have 
$r_{\beta}wT_{\xi} = w r_{\gamma} T_{\xi}$ 
and 
$wr_{\gamma} \in W^J$.\setcounter{enumi}{16}
\item\label{Prop2.5_Q} %[(Q)]
$\chi = 1$ 
and 
$\ell(\lfloor w r_{\gamma} \rfloor) = \ell(w) + 1 - 2( \gamma^{\vee} , \rho - \rho_J )$; 
in this case, we have
$r_{\beta}wT_{\xi} = \lfloor w r_{\gamma} \rfloor T_{\xi + \gamma^{\vee}}$.
\end{enumerate}
\end{prop}

\begin{rema}\ \\[-1.2em]
\begin{enumerate}[label=(\arabic*)]
\item\label{Rema2.6_1}
If 
$wT_{\xi} \xrightarrow{\ \beta \ } \Pi^J(r_{\beta}wT_{\xi})$
in 
$\SiB ^J$, 
then 
$r_{\beta}wT_{\xi} = \Pi^J (r_{\beta}wT_{\xi}) \in (W^J)_{\af}$
(see~\cite[Appendix~A]{INS16}). 
\item\label{Rema2.6_2}
The condition~\ref{Prop2.5_B} (\resp~~\ref{Prop2.5_Q}) for 
$w \in W^J$ 
and 
$\gamma \in \Delta^+ \setminus \Delta^+_J$ 
in Proposition~\ref{prop:phi_J} corresponds to the existence of the 
Bruhat edge (\resp quantum edge) 
$w \xrightarrow{\ \gamma \ } \lfloor w r_{\gamma} \rfloor$
in the (parabolic) quantum Bruhat graph for 
$W^J$
(see~\cite[\S~4]{LNSSS15}).
\end{enumerate}
\end{rema}

\subsection{Semi-infinite Lakshmibai--Seshadri paths}

In this subsection, 
we give a brief exposition of the 
$\U$-crystal of semi-infinite Lakshmibai--Seshadri paths; 
see~\cite{INS16} for more details.

Let 
$\lambda \in P^+$
and set 
$J_{\lambda} = \{ j \in I \mid \langle h_j , \lambda \rangle = 0 \}$. 
For a rational number 
$0 < a \leq 1$, 
define 
$\SiB (\lambda;a)$
to be the subgraph of 
$\SiB ^{J_{\lambda}}$ 
with the same vertex set 
but having only the edges of the form 
\begin{align}
x \xrightarrow{\ \beta \ } y
\ \ \text{with} \ \ 
a ( \beta^{\vee} , x\lambda ) \in \BZ ;
\end{align}
note that 
$\SiB (\lambda ; 1) = \SiB ^{J_{\lambda}}$. 
A semi-infinite Lakshmibai--Seshadri path of shape $\lambda$ is, by definition, 
a pair 
$(\bm{x}; \bm{a})$
of a decreasing sequence 
$\bm{x} : x_1 \succeq x_2 \succeq \cdots \succeq x_l$ 
of elements in $(W^{J_{\lambda}})_{\af}$
and an increasing sequence
$\bm{a} : 0 = a_0 < a_1 < \cdots < a_l = 1$ 
of rational numbers such that 
there exists a directed path from
$x_{u+1}$ to $x_u$ in $\SiB (\lambda;a_u)$ 
for each $u=1,2,\ldots ,l-1$. 
Let 
$\mathbb{B}^{\si}(\lambda)$ 
denote the set of semi-infinite Lakshmibai--Seshadri paths of shape $\lambda$. 

Following~\cite[\S~3.1]{INS16}, 
we equip the set 
$\mathbb{B}^{\si}(\lambda)$ 
with a $\U$-crystal structure.
For 
$\eta = (x_1 , \ldots , x_l ; a_0 , a_1 , \ldots , a_l) \in \mathbb{B}^{\si}(\lambda)$, 
define the map 
$\bar{\eta} : [0,1] \rightarrow \BR \otimes_{\BZ} P_{\af}$ 
by
\begin{align} \label{eq:path}
\bar{\eta}(t) 
=
\sum_{p=1}^{u-1} (a_p - a_{p-1}) x_p \lambda + (t - a_{u-1}) x_u \lambda 
\ \ 
\text{for} \ 
t \in [a_{u-1} , a_u], \ 
1 \leq u \leq l.
\end{align}
Define 
$\wt : \mathbb{B}^{\si}(\lambda) \rightarrow P_{\af}$
by
$\wt (\eta) = \bar{\eta}(1) \in P_{\af}$.
Set
\begin{align} \label{eq:path-h-m}
h_i^{\eta}(t) 
= 
\langle h_i , \bar{\eta}(t) \rangle 
\ \ \text{for}\ 
t \in [0,1] , &&
m_i^{\eta} 
= 
\min \{ h_i^{\eta}(t) \mid t \in [0,1] \} .
\end{align}

We define 
$e_i \eta$, $i \in I_{\af}$, 
as follows: 
if $m_i^{\eta} = 0$, 
then we set $e_i \eta = \bm{0}$.
If $m_i^{\eta} \leq -1$, then we set
\begin{align} \label{eq:e_i-t}
\begin{cases}
t_1 = \min \{ t \in [0,1] \mid h_i^{\eta}(t) = m_i^{\eta} \} , \\%[1mm]
t_0 = \max \{ t \in [0,t_1] \mid h_i^{\eta}(t) = m_i^{\eta} + 1 \} .
\end{cases}
\end{align}
Let 
$1 \leq p \leq q \leq l$
be such that 
$a_{p-1} \leq t_0 < a_p$ and $t_1 = a_q$. 
Then we define 
\begin{align} \label{eq:e_i}
\begin{split}
e_i \eta 
= 
(x_1 , \ldots , x_p , r_i x_p , \ldots , 
&\, r_i x_q , x_{q+1} , \ldots , x_l ; \\
&a_0 , \ldots , a_{p-1} , t_0 , a_p , \ldots , a_q = t_1 , \ldots , a_l ) ;
\end{split}
\end{align}
if $t_0 = a_{p-1}$, 
then we drop $x_p$ and $a_{p-1}$, 
and if $r_j x_q = x_{q+1}$, 
then we drop $x_{q+1}$ and $a_q = t_1$.

Next, we define $f_i \eta$, $i \in I_{\af}$, as follows: 
if $m_i^{\eta} = h_i^{\eta}(1)$, then we set $f_i \eta = \bm{0}$. 
If $h_i^{\eta}(1) - m_i^{\eta} \geq 1$, then we set
\begin{align} \label{eq:f_i-t}
\begin{cases}
t_0 = \max \{ t \in [0,1] \mid h_i^{\eta}(t) = m_i^{\eta} \} , \\%[1mm]
t_1 = \min \{ t \in [t_0,1] \mid h_i^{\eta}(t) = m_i^{\eta} + 1 \} .
\end{cases}
\end{align}
Let 
$1 \leq p \leq q \leq l-1$
be such that
$t_0 = a_p$ and $a_q < t_1 \leq a_{q+1}$. 
Then we define
\begin{align}
\begin{split}
f_i \eta 
= 
(x_1 , \ldots , x_p , r_i x_{p+1} , \ldots , 
& \, r_i x_{q+1} , x_{q+1} , \ldots , x_l ; \\ 
&a_0 , \ldots , a_p = t_0 , \ldots , a_q , t_1 , a_{q+1} , \ldots , a_l ) ;
\end{split}
\end{align}
if $t_1 = a_{q+1}$, 
then we drop $x_{q+1}$ and $a_{q+1}$, 
and if $x_p = r_i x_{p+1}$, 
then we drop $x_p$ and $a_p = t_0$.

For $\eta \in \mathbb{B}^{\si}(\lambda)$ and $i \in I_{\af}$, define
\begin{align} \label{eq:path-eps-phi}
\begin{cases}
\varepsilon_i(\eta) = -m_i^{\eta} , \\%[1mm]
\varphi_i(\eta) = h_i^{\eta}(1) - m_i^{\eta} .
\end{cases}
\end{align}

Now we assume that 
$\lambda = \sum_{i \in I} m_i \varpi_i$, 
with 
$m_i \in \BZ_{\geq 0}$, 
$i \in I$. 
Set 
$J_{\lambda}^c = I \setminus J_{\lambda} = \{ i \in I \mid m_i > 0 \}$. 
Following~\cite[Equation~(7.2.2)]{INS16}, 
we define an element 
$\eta_{\bmrho{\rho}} \in \mathbb{B}^{\si}(\lambda)$ 
of weight 
$\lambda + \wt (\bmrho{\rho})$
for each 
$\bmrho{\rho} = \left( \rho^{(i)} \right)_{i \in I} \in \Par (\lambda)$, 
with 
$\rho^{(i)} = \left(\rho^{(i)}_1 \geq \rho^{(i)}_2 \geq \cdots \geq \rho^{(i)}_{m_i-1}\right)$
(see~\eqref{eq:Par}). 
Let 
$s$ 
be the least common multiple of 
$\{ m_i \mid i \in J_{\lambda}^c \}$. 
Let 
$c_i (\xi) \in \BZ$ 
denote the coefficient of 
$\alpha_i^{\vee}$ 
in 
$\xi \in Q^{\vee}$.
For 
$\xi , \zeta \in Q^{\vee}$, 
write 
$\xi \succeq \zeta$ 
if 
$\xi - \zeta \in \sum_{i \in I} \BZ_{\geq 0} \alpha_i^{\vee}$, 
and write 
$\xi \succ \zeta$
if 
$\xi \succeq \zeta$ 
and 
$\xi \neq \zeta$. 
Let 
$\zeta_1 , \ldots , \zeta_s \in Q^{\vee}$
be such that 
\begin{enumerate}[label=(\roman*)]
\item
$c_i(\zeta_t) = \rho^{(i)}_u$ 
if $i \in J_{\lambda}^c$ 
and $\dfrac{s(u-1)}{m_i} < t \leq \dfrac{su}{m_i}$, and 
\item
$c_j(\zeta_t) = 0$ 
for all $j \in J_{\lambda}$ and $1 \leq t \leq s$; 
\end{enumerate}
note that 
$\zeta_1 \succeq \cdots \succeq \zeta_s$
and 
$\zeta_s = 0$. 
Assume that
\begin{align}
\zeta_1 = \cdots = \zeta_{s_1} \succ \zeta_{s_1 + 1} = \cdots = \zeta_{s_2} \succ \cdots \cdots 
\succ \zeta_{s_{k-1} + 1} = \cdots = \zeta_{s_k}, 
\end{align}
where 
$1 \leq s_1 < \cdots < s_{k-1} < s_k = s$.
Set 
\begin{align}
\eta_{\bmrho{\rho}} 
=
\left( 
T_{\zeta_{s_1}} , \ldots , T_{\zeta_{s_{k-1}}} , e ; \, 
0 , \dfrac{s_1}{s} , \ldots , \dfrac{s_{k-1}}{s} , 1
\right) .
\end{align}

\begin{theo}[{\cite[Theorems~3.1.5 and~3.2.1]{INS16}}] \label{thm:SiLS-isom}
Let $\lambda \in P^+$.
\begin{enumerate}[label=(\arabic*)]
\item\label{Theo2.7_1}
The set
$\mathbb{B}^{\si}(\lambda)$
equipped with the maps 
$\wt $, 
$e_i , f_i$, $i \in I_{\af}$, and
$\varepsilon_i , \varphi_i$, $i \in I_{\af}$, 
defined above, is a $\U$-crystal. 
\item\label{Theo2.7_2}
There exists a unique 
isomorphism of $\U$-crystals from 
$\mathcal{B}(\lambda)$
to 
$\mathbb{B}^{\si}(\lambda)$
sending 
$S_{\bmrho{\rho}}^- u_{\lambda} \mod q_s \mathcal{L}(\lambda)$ 
to 
$\eta_{\bmrho{\rho}}$
for every 
$\bmrho{\rho} \in \Par (\lambda)$. 
\end{enumerate}
\end{theo}

For 
$\lambda_1 , \ldots , \lambda_N \in P^+$, 
let 
$\mathbb{B}^{\si}(\lambda_1) * \cdots * \mathbb{B}^{\si}(\lambda_N)$
be the set of symbols
$\eta_1 * \cdots * \eta_N$, 
with 
$\eta_{\nu} \in \mathbb{B}^{\si}(\lambda_{\nu})$, 
$1 \leq \nu \leq N$. 
We define a $\U$-crystal structure on 
$\mathbb{B}^{\si}(\lambda_1) * \cdots * \mathbb{B}^{\si}(\lambda_N)$
in a way similar to the above. For 
$\eta = \eta_1 * \cdots * \eta_N \in \mathbb{B}^{\si}(\lambda_1) * \cdots * \mathbb{B}^{\si}(\lambda_N)$, 
define 
$\bar{\eta} : [0,1] \rightarrow \BR \otimes_{\BZ} P_{\af}$ 
by
\begin{align} \label{eq:con-path}
\bar{\eta} (t) = 
\sum_{\nu = 1}^{\mu - 1} \bar{\eta}_{\nu}(1) + \bar{\eta}_{\mu}(Nt - \mu + 1) 
\ \ 
\text{for} \ 
\dfrac{\mu - 1}{N} \leq t \leq \dfrac{\mu}{N}, \ 
1 \leq \mu \leq N,
\end{align}
where each 
$\bar{\eta}_{\nu} : [0,1] \rightarrow \BR \otimes_{\BZ} P_{\af}$, $1 \leq \nu \leq N$, 
is defined by~\eqref{eq:path}. 
Define 
$\wt (\eta) = \bar{\eta}(1)$. 
By the same way as in~\eqref{eq:path-h-m}, 
we define 
$h_i^{\eta}(t)$ 
and 
$m_i^{\eta}$
for 
$\eta = \eta_1 * \cdots * \eta_N$ 
by using~\eqref{eq:con-path}. 
We define 
$e_i \eta \ (\resp f_i \eta) \in \mathbb{B}^{\si}(\lambda_1) * \cdots * \mathbb{B}^{\si}(\lambda_N) \sqcup \{ \bm{0} \}$ 
as follows: 
if $m_i^{\eta} = 0$ (\resp $m_i^{\eta} = h_i^{\eta} (1) $), then we set $e_i \eta = \bm{0}$ (\resp $f_i \eta = \bm{0}$). 
Assume that 
$m_i^{\eta} \leq -1$ 
(\resp $h_i^{\eta} (1) - m_i^{\eta} \geq 1$), 
and let 
$0 \leq t_0 < t_1 \leq 1$ 
be as in~\eqref{eq:e_i-t} (\resp~\eqref{eq:f_i-t}). 
We see that there exists $1 \leq \nu \leq N$ such that 
$\dfrac{\nu -1}{N} \leq t_0 < t_1 \leq \dfrac{\nu}{N}$. 
Then we set 
$e_i \eta = \eta_1 * \cdots * \eta_{\nu -1} * e_i \eta_{\nu} * \eta_{\nu + 1} * \cdots * \eta_N$ 
(\resp $f_i \eta = \eta_1 * \cdots * \eta_{\nu -1} * f_i \eta_{\nu} * \eta_{\nu + 1} * \cdots * \eta_N$).
We define the functions 
$\varepsilon_i , \varphi_i$ 
as in~\eqref{eq:path-eps-phi}. 
The proof of the next proposition is straightforward.

\begin{prop} \label{prop:conc=tensor}
Let 
$\lambda_1 , \ldots , \lambda_N \in P^+$. 
The map 
$\mathbb{B}^{\si}(\lambda_1) \otimes \cdots \otimes \mathbb{B}^{\si}(\lambda_N)
\rightarrow 
\mathbb{B}^{\si}(\lambda_1) * \cdots * \mathbb{B}^{\si}(\lambda_N)$, 
$\eta_1 \otimes \cdots \otimes \eta_N \mapsto \eta_1 * \cdots * \eta_N$, 
is an isomorphism of $\U$-crystals.
\end{prop}

\section{Standard monomial theory for semi-infinite Lakshmibai--Seshadri paths} \label{Section:SMT}

\subsection{Strict embedding \texorpdfstring{$\Phi_{\lambda \mid q=0}^{\LT }$}{Phi lambda|q=0 LT}} \label{Subsection:LT}

Let 
$\lambda = \sum_{i \in I} m_i \varpi_i \in P^+$, 
with 
$m_i \in \BZ_{\geq 0}$, 
$i \in I$. 
Recall the automorphisms 
$z_{i,\nu}$, 
$i \in I$, $1 \leq \nu \leq m_i$, 
of the $\U'$-crystal
$\Breve{\mathcal{B}}(\lambda)$
(see \S~\ref{Subsection:extremal}). 
For 
$\bmrho{\rho} = (\rho^{(i)})_{i \in I} \in \Par (\lambda)$, 
with 
$\rho^{(i)} = (\rho^{(i)}_1 \geq \rho^{(i)}_2 \geq \cdots \geq \rho^{(i)}_{m_i -1} \geq 0)$, 
$i \in I$, 
define the automorphism 
$z^{-\bmrho{\rho}}$ 
of the $\U'$-crystal 
$\Breve{\mathcal{B}}(\lambda)$
by
\begin{align}
z^{-\bmrho{\rho}}
= 
\prod_{i \in I}
z_{i,1}^{- \rho^{(i)}_1} 
z_{i,2}^{- \rho^{(i)}_2}
\cdots 
z_{i,m_i-1}^{- \rho^{(i)}_{m_i-1}} .
\end{align}
Define the map
$\Phi_{\lambda \mid q=0}^{\LT } : 
\mathcal{B}(\lambda) \rightarrow 
\Breve{\mathcal{B}}(\lambda)$
by 
\begin{align} \label{eq:LT} 
g_1 g_2 \cdots g_l S_{\bmrho{\rho}}^- u_{\lambda} 
\mod q_s \mathcal{L}(\lambda) 
\ \longmapsto \ 
g_1 g_2 \cdots g_l z^{-\bmrho{\rho}} \Tilde{u}_{\lambda} 
\mod q_s \Breve{\mathcal{L}}(\lambda) ,
\end{align}
where 
$g_k \in \{ e_i , f_i \mid i \in I_{\af} \}$, 
$1 \leq k \leq l$, 
$l \in \BZ_{\geq 0}$, and 
$\bmrho{\rho} \in \Par (\lambda)$ 
(\cf~\eqref{eq:Phi_q=0}). 
Set 
$\Breve{\mathcal{B}}_{\bmrho{\rho}}(\lambda)
=
z^{-\bmrho{\rho}} \left( \Breve{\mathcal{B}}_0(\lambda) \right)
\subset \Breve{\mathcal{B}}(\lambda)$; 
note that
$\Breve{\mathcal{B}}_{\bmrho{\rho}}(\lambda)$
is a connected component of 
$\Breve{\mathcal{B}}(\lambda)$, 
and is isomorphic to
$\Breve{\mathcal{B}}_0(\lambda)$
as a $\U'$-crystal.

\begin{lemm} \label{lem:Phi_LT}
The map 
$\Phi_{\lambda \mid q=0}^{\LT }$ 
is well-defined, and is a strict embedding of $\U$-crystals.
\end{lemm}

\begin{proof}
It suffices to show that the map
$\Phi_{\lambda \mid q=0}^{\LT }$ 
induces an isomorphism of $\U$-crystals from
$\mathcal{B}_{\bmrho{\rho}}(\lambda)$
to
$\Breve{\mathcal{B}}_{\bmrho{\rho}}(\lambda)$
for every 
$\bmrho{\rho} \in \Par (\lambda)$.
We know that the maps
$\mathcal{B}_{\bmrho{\rho}}(\lambda) \rightarrow \mathcal{B}_0(\lambda)$ 
in Theorem~\ref{thm:BN}~\ref{Theo2.1_3}, 
$\Phi_{\lambda \mid q=0} : 
\mathcal{B}_0(\lambda)
\rightarrow 
\Breve{\mathcal{B}}_0(\lambda)$
in Theorem~\ref{thm:BN}~\ref{Theo2.1_2}, and 
$z^{-\bmrho{\rho}} : 
\Breve{\mathcal{B}}_0(\lambda)
\rightarrow 
\Breve{\mathcal{B}}_{\bmrho{\rho}}(\lambda)$
are isomorphisms of $\U'$-crystals. 
We check at once that the composition of these maps is describe by~\eqref{eq:LT}, 
which proves that 
$\Phi_{\lambda \mid q=0}^{\LT }$ 
is well-defined and induces an isomorphism of $\U'$-crystals from 
$\mathcal{B}_{\bmrho{\rho}}(\lambda)$
to
$\Breve{\mathcal{B}}_{\bmrho{\rho}}(\lambda)$. 
Since 
$\wt (S_{\bmrho{\rho}}^- u_{\lambda}) 
= 
\lambda + \wt (\bmrho{\rho})
=
\wt (z^{-\bmrho{\rho}} \Tilde{u}_{\lambda})$, 
$\Phi_{\lambda \mid q=0}^{\LT }$ 
is a morphism of $\U$-crystals. 
\end{proof}

\begin{rema} \label{rem:LT}
If we think of the Schur polynomials
as the generating functions of the 
weights of Young tableaux (see~\cite[Page~3]{F}), 
then the term
$z^{-\bmrho{\rho}}$
in 
$s_{\bmrho{\rho}}(z^{-1})$ 
(see~\eqref{eq:Schur})
corresponds to the tuple 
of the Littlewood--Richardson tableaux 
of shapes $\rho^{(i)}$, $i \in I$
(see~\cite[\S~5.2]{F}), 
and the coefficient of 
$z^{-\bmrho{\rho}}$
in 
$s_{\bmrho{\rho}}(z^{-1})$
is $1$. 
\end{rema}

\subsection{Characterization of the image of \texorpdfstring{$\Phi_{\lambda \mid q=0}^{\LT }$}{Phi lambda|q=0 LT}}

In this subsection, we give a characterization of the image of the map 
$\Phi_{\lambda \mid q=0}^{\LT }$
in terms of semi-infinite Bruhat order 
via semi-infinite Lakshmibai--Seshadri paths. 

Recall the notation 
$J_{\lambda} = \{ i \in I \mid \langle h_i , \lambda \rangle = 0 \}$, 
$\lambda \in P^+$. 

\begin{defi}
Let 
$\lambda_{\nu} \in P^+$
and 
$\eta^{(\nu)} = \left( x^{(\nu)}_1 , \ldots , x^{(\nu)}_{l_{\nu}} ;\, 
\bm{a}^{(\nu)} \right) \in \mathbb{B}^{\si}(\lambda_{\nu})$, 
$1 \leq \nu \leq N$. 
We say that there exists a defining chain for
$\bigotimes_{\nu = 1}^N \eta^{(\nu)} \in \bigotimes_{\nu = 1}^N \mathbb{B}^{\si}(\lambda_{\nu})$ 
if there exists 
$\Tilde{x}^{(\nu)}_s \in W_{\af}$, 
$1 \leq s \leq l_{\nu}$, $1 \leq \nu \leq N$, 
such that 
\begin{enumerate}[label=(DC\arabic*)]
\item
$\Pi^{J_{\lambda_{\nu}}}(\Tilde{x}^{(\nu)}_s) = x^{(\nu)}_s$ 
for all $1 \leq s \leq l_{\nu}$, $1 \leq \nu \leq N$, 
\item
$\Tilde{x}^{(\nu)}_p \succeq \Tilde{x}^{(\nu)}_q$ 
for all
$1 \leq p \leq q \leq l_{\nu}$, 
$1 \leq \nu \leq N$, 
and 
\item
$\Tilde{x}^{(\nu)}_{l_{\nu}} \succeq \Tilde{x}^{(\nu + 1)}_1$ 
for all $1 \leq \nu < N$, 
\end{enumerate}
where $\succeq$ 
denotes the semi-infinite Bruhat order 
on $W_{\af}$ defined by using 
$\SiB ^{\emptyset}$
(see \S~\ref{Subsection:SiBO}). 
The tuple
$( \Tilde{x}^{(\nu)}_s )_{1 \leq s \leq l_{\nu}, 1 \leq \nu \leq N}$
above is called a defining chain for 
$\bigotimes_{\nu = 1}^N \eta^{(\nu)}$.
\end{defi}

Let 
$\lambda = \sum_{i \in I} m_i \varpi_i \in P^+$. 
Recall the notation
$J_{\lambda}^c = I \setminus J_{\lambda}$. 
Set
$\Breve{\mathbb{B}}^{\si}(\lambda) = \bigotimes_{i \in J_{\lambda}^c} \mathbb{B}^{\si}(\varpi_i)^{\otimes m_i}$.
We know from Theorem~\ref{thm:SiLS-isom} that there exists an isomorphism 
\begin{align}
\Psi_{\lambda} : \Breve{\mathcal{B}}(\lambda) \rightarrow \Breve{\mathbb{B}}^{\si}(\lambda)
\end{align}
of $\U$-crystals 
defined as the tensor product of the isomorphisms 
$\mathcal{B}(\varpi_i) \rightarrow \mathbb{B}^{\si}(\varpi_i)$
sending 
$u_{\varpi_i}$ 
to 
$(e;0,1)$, 
$i \in I$. 
Write
\begin{align}
\Psi_{\lambda}(b) 
= 
\bigotimes_{i \in J_{\lambda}^c} \Psi^{(i)}_{\lambda}(b) \in \Breve{\mathbb{B}}^{\si}(\lambda), 
\ \text{where} \ 
\Psi^{(i)}_{\lambda}(b) \in \mathbb{B}^{\si}(\varpi_i)^{\otimes m_i}, \ 
i \in J_{\lambda}^c .
\end{align}
Set 
\begin{align}
\Breve{\mathbb{S}}^{\si}(\lambda) 
= 
\left\{ \eta \in \Breve{\mathbb{B}}^{\si}(\lambda) 
\ \vline \ 
\text{there exists a defining chain for} \ \eta \right\} .
\end{align}

\goodbreak
\begin{theo} \label{thm:SMT}
Let $\lambda \in P^+$. 
For $b \in \Breve{\mathcal{B}}(\lambda)$, 
the following conditions are equivalent:
\begin{enumerate}[label=(\arabic*)]
\item\label{Theo3.4_1}
$b \in \Phi_{\lambda \mid q=0}^{\LT }(\mathcal{B}(\lambda))$.
\item\label{Theo3.4_2}
$\Psi_{\lambda}(b) \in \Breve{\mathbb{S}}^{\si}(\lambda)$. 
\item\label{Theo3.4_3}
$\Psi_{\lambda}^{(i)}(b) \in \Breve{\mathbb{S}}^{\si}(m_i \varpi_i)$
for every $i \in J_{\lambda}^c$.
\end{enumerate}
In particular, 
$\Breve{\mathbb{S}}^{\si}(\lambda) = \bigotimes_{i \in J_{\lambda}^c}\Breve{\mathbb{S}}^{\si}(m_i \varpi_i)$ 
(\cf~\cite[Remark~4.17]{BN04}; see also~\cite[Conjecture~13.1~(iii)]{K02}), 
$\Breve{\mathbb{S}}^{\si}(\lambda)$ 
is stable under the Kashiwara operators, 
and the map 
$\Psi_{\lambda} \circ \Phi_{\lambda \mid q=0}^{\LT } : \mathcal{B}(\lambda) \rightarrow \Breve{\mathbb{S}}^{\si}(\lambda)$ 
is an isomorphism of $\U$-crystals.
\end{theo}

\goodbreak
\begin{rema}\ \\[-1.2em]
\begin{enumerate}[label=(\arabic*)]
\item\label{Rema3.5_1}
Our argument in the proof of Theorem~\ref{thm:SMT} in \S~\ref{Subsection:pf-SMT} 
does not imply~\cite[Conjecture~13.1~(iii)]{K02} 
since~\cite[Remark~4.17]{BN04}
is used in the proof of 
Theorem~\ref{thm:SiLS-isom}
(\cite[Theorems~3.1.5 and~3.2.1]{INS16}). 
\item\label{Rema3.5_2}
Similar result to Theorem~\ref{thm:SMT} is obtained in~\cite[Theorem~3.1]{KNS}, 
where they proved that, 
for any $\lambda , \mu \in P^+$, 
the subset 
$\mathbb{S}^{\si}(\lambda + \mu)$
of elements in 
$\mathbb{B}^{\si}(\lambda) \otimes \mathbb{B}^{\si}(\mu)$
having a defining chain is stable under the Kashiwara operators, 
and is isomorphic to 
$\mathbb{B}^{\si}(\lambda + \mu)$
as a $\U$-crystal. 
But the proof is slightly different from ours.
The main task in the proof of~\cite[Theorem~3.1]{KNS}
is to construct an isomorphism of $\U$-crystals between 
$\mathbb{S}^{\si}(\lambda + \mu)$
and 
$\mathbb{B}^{\si}(\lambda + \mu)$.
This is achieved by giving an explicit parametrization of the 
connected components of $\mathbb{S}^{\si}(\lambda + \mu)$.
In contrast of this, our argument starts from a specific choice of a map (see~\eqref{eq:LT}), 
and aims to give an explicit description of the image of this map; 
in fact, there are infinitely many strict embeddings of $\U$-crystals from
$\mathcal{B}(\lambda)$
to
$\Breve{\mathcal{B}}(\lambda)$, 
in general.
\end{enumerate}
\end{rema}

\subsection{Proof of Theorem~\ref{thm:SMT}} \label{Subsection:pf-SMT}
This subsection is devoted to the proof of Theorem~\ref{thm:SMT}.

We see from~\cite[Theorem~5.17]{K02} 
that, for each $i \in I$, there exists a strict surjective morphism of $\U'$-crystals from 
$\mathcal{B}(\varpi_i)$
to the crystal basis of a finite-dimensional $\U'$-module
$W(\varpi_i)$ (see \S~\ref{Subsection:extremal}). 
Hence, the next lemma follows from~\cite[Lemmas~1.5--1.6]{AK}. 

\begin{lemm} \label{lem:extremal-vectors}
Let 
$i_1 , \ldots , i_N \in I$. 
\begin{enumerate}[label=(\arabic*)]
\item\label{Theo3.6_1}
Any connected component of 
$\bigotimes_{\nu = 1}^N \mathcal{B}(\varpi_{i_{\nu}})$ 
contains an extremal element.
\item\label{Theo3.6_2}
If 
$b = \bigotimes_{\nu = 1}^N b^{(\nu)} \in \bigotimes_{\nu = 1}^N \mathcal{B}(\varpi_{i_{\nu}})$
is an extremal element, then 
$S_x b = \bigotimes_{\nu = 1}^N S_x b^{(\nu)}$ 
for all 
$x \in W_{\af}$. 
\item\label{Theo3.6_3}
$\bigotimes_{\nu = 1}^N b^{(\nu)} \in \bigotimes_{\nu = 1}^N \mathcal{B}(\varpi_{i_{\nu}})$
is an extremal element if and only if there exist
$w \in W$ 
and 
$\xi_1 , \ldots , \xi_N \in Q^{\vee}$ 
such that 
$b^{(\nu)} = S_{wt_{\xi_{\nu}}} u_{\varpi_{i_{\nu}}}$
for 
$1 \leq \nu \leq N$. 
\end{enumerate}
\end{lemm}

By Theorem~\ref{thm:SiLS-isom}, 
the $W_{\af}$-action on $\mathcal{B}(\lambda)$ (see~\eqref{eq:W-action})
induces a $W_{\af}$-action on $\mathbb{B}^{\si}(\lambda)$.
% as follows: 
%$S_x (x_1 , \ldots , x_l ; \bm{a}) = (\Pi^{J_{\lambda}}(xx_1) , \ldots , \Pi^{J_{\lambda}}(xx_l) ; \bm{a})$ 
%for 
%$x \in W_{\af}$
%and 
%$(x_1 , \ldots , x_l ; \bm{a}) \in \mathbb{B}^{\si}(\lambda)$.

The next lemma follows from 
Theorem~\ref{thm:SiLS-isom}~\ref{Theo2.7_2} and 
Lemmas~\ref{lem:Pi}--\ref{lem:T_xi} and~\ref{lem:extremal-vectors}.

\begin{lemm} \label{lem:extremal-paths}
Let 
$i_1 , \ldots , i_N \in I$. 
\begin{enumerate}[label=(\arabic*)]
\item\label{Lemm3.7_1}
Any connected component of 
$\bigotimes_{\nu = 1}^N \mathbb{B}^{\si}(\varpi_{i_{\nu}})$ 
contains an extremal element.
\item\label{Lemm3.7_2}
If 
$\eta = \bigotimes_{\nu = 1}^N \eta^{(\nu)} \in \bigotimes_{\nu = 1}^N \mathbb{B}^{\si}(\varpi_{i_{\nu}})$
is an extremal element, then 
$S_x \eta = \bigotimes_{\nu = 1}^N S_x \eta^{(\nu)}$
for all 
$x \in W_{\af}$. 
\item\label{Lemm3.7_3}
$\bigotimes_{\nu = 1}^N \eta^{(\nu)} \in \bigotimes_{\nu = 1}^N \mathbb{B}^{\si}(\varpi_{i_{\nu}})$
is an extremal element if and only if there exist
$w \in W$ 
and 
$\sigma_{\nu} \in \BZ$, $1 \leq \nu \leq N$, 
such that 
$\eta^{(\nu)} 
= 
S_w \left( T^{I \setminus \{ i_{\nu} \}}_{\sigma_{\nu} \alpha_{i_{\nu}}^{\vee}} ; 0,1 \right) \in \mathbb{B}^{\si}(\varpi_{i_{\nu}})$
for 
$1 \leq \nu \leq N$. 
\end{enumerate}
\end{lemm}

We also denote by
$z_i$ 
the automorphism, as a $\U'$-crystal, of 
$\mathbb{B}^{\si}(\varpi_i)$ 
corresponding to the automorphism 
$z_i$ 
of 
$\mathcal{B}(\varpi_i)$.
Recall that 
$c_i(\xi)$ 
denotes the coefficient of 
$\alpha_i^{\vee}$ 
in 
$\xi \in Q^{\vee}$.

\begin{lemm} \label{lem:SiLS_W-action}
Let 
$\lambda = \sum_{i \in I} m_i \varpi_i \in P^+$. 
\begin{enumerate}[label=(\arabic*)]
\item\label{Lemm3.8_1}
For each $i \in I$, 
we have
$z_i^k(e;0,1) = \left( T^{I \setminus \{ i \}}_{-k \alpha_i^{\vee}} ; 0,1 \right)$
in 
$\mathbb{B}^{\si}(\varpi_i)$
for all 
$k \in \BZ$. 
In particular, 
for every 
$\bmrho{\rho} = (\rho^{(i)})_{i \in I} \in \Par (\lambda)$, 
with 
$\rho^{(i)} = (\rho^{(i)}_1 \geq \cdots \geq \rho^{(i)}_{m_i-1} \geq 0)$, 
$i \in I$, 
we have
\begin{align}
\Psi_{\lambda}^{(i)}(z^{-\bmrho{\rho}} \Tilde{u}_{\lambda}) 
=
\left( T_{\rho^{(i)}_1 \alpha_i^{\vee}}^{I \setminus \{ i \}} ; 0,1\right)
\otimes \cdots \otimes
\left( T_{\rho^{(i)}_{m_i-1} \alpha_i^{\vee}}^{I \setminus \{ i \}} ; 0,1\right)
\otimes
(e;0,1)
\end{align}
in 
$\mathbb{B}^{\si}(\varpi_i)^{\otimes m_i}$
for each 
$i \in I$. 
\item\label{Lemm3.8_2}
Any connected component of 
$\Breve{\mathbb{B}}^{\si}(\lambda)$
contains an extremal element of the form 
\begin{align}
\bigotimes_{i \in J_{\lambda}^c} \bigotimes_{\nu = 1}^{m_i} \left(T^{I \setminus \{ i \}}_{\rho^{(i)}_{\nu} \alpha_i^{\vee}};0,1 \right),
\end{align}
where 
$\rho^{(i)}_{\nu} \in \BZ$, 
$1 \leq \nu \leq m_i$, 
$i \in J_{\lambda}^c$, 
and 
$\rho^{(i)}_{m_i} = 0$
for all 
$i \in J_{\lambda}^c$.
\end{enumerate}
\end{lemm}

\begin{proof}
\ref{Lemm3.8_1}: 
Since 
$z_i^k(e;0,1)$
is an extremal element of weight $\varpi_i + k\delta$, 
there exists 
$x \in (W^{I \setminus \{ i \}})_{\af}$
such that 
$z_i^k(e;0,1) = S_x (e;0,1) = (x;0,1)$ 
by~\cite[Proposition~5.4~(i)]{K02}. 
If we write 
$x = wT^{I \setminus \{ i \}}_{\xi}$, 
$w \in W$, 
$\xi \in Q^{\vee}$, 
then 
$\wt (x;0,1) 
= w\varpi_i - (\xi,\varpi_i)\delta = w \varpi_i - c_i(\xi) \delta$, 
which implies that 
$w = e$
and 
$c_i(\xi) = -k$. 
By Lemma~\ref{lem:T_xi}~\ref{Lemm2.4_2}, 
we have 
$T^{I \setminus \{ i \}}_{\xi} = T^{I \setminus \{ i \}}_{-k\alpha_i^{\vee}}$, 
which proves that 
$z_i^k(e;0,1) = \left( T^{I \setminus \{ i \}}_{-k\alpha_i^{\vee}} ; 0,1 \right)$. 
 
\ref{Lemm3.8_2}: 
By Lemma~\ref{lem:extremal-paths}, 
any connected component $C$ of 
$\Breve{\mathbb{B}}^{\si}(\lambda)$
contains an extremal element of the form 
$\eta = \bigotimes_{i \in J_{\lambda}^c} 
\bigotimes_{\nu=1}^{m_i} \left(T^{I \setminus \{ i \}}_{\sigma^{(i)}_{\nu} \alpha_i^{\vee}} ;0,1 \right)$, 
with 
$\sigma^{(i)}_{\nu} \in \BZ$, $1 \leq \nu \leq m_i$, $i \in J_{\lambda}^c$. 
Set 
$\rho^{(i)}_{\nu} = \sigma^{(i)}_{\nu} - \sigma^{(i)}_{m_i}$ 
and 
$\xi = -\sum_{i \in J_{\lambda}^c} \sigma^{(i)}_{m_i} \alpha_i^{\vee}$. 
By Lemmas~\ref{lem:T_xi}~\ref{Lemm2.4_2} and~\ref{lem:extremal-paths}~\ref{Lemm3.7_2}, 
\begin{align}
C \ni \, 
S_{t_{\xi}} \eta 
= 
\bigotimes_{i \in J_{\lambda}^c} \bigotimes_{\nu = 1}^{m_i} 
S_{t_{\xi}} \left(T^{I \setminus \{ i \}}_{\sigma^{(i)}_{\nu} \alpha_i^{\vee}} ;0,1 \right)
=
\bigotimes_{i \in J_{\lambda}^c} \bigotimes_{\nu = 1}^{m_i} 
\left( T^{I \setminus \{ i \}}_{\rho^{(i)}_{\nu} \alpha_i^{\vee}} ; 0,1 \right) , 
\end{align}
which is the desired conclusion.
\end{proof}

For 
$J \subset I$, 
set
$J^c = I \setminus J$. 
Let 
$[ \, \cdot \, ]_J : Q^{\vee} = Q^{\vee}_J \oplus Q^{\vee}_{J^c} \rightarrow Q_J^{\vee}$
be the projection. 
Recall that we write
$\xi \succeq \zeta$
for 
$\xi , \zeta \in Q^{\vee}$ 
if 
$\xi - \zeta \in \sum_{i \in I} \BZ_{\geq 0} \alpha_i^{\vee}$.

\begin{lemm} \label{lem:easy}\ \\[-1.2em]
\begin{enumerate}[label=(\arabic*)]
\item\label{Lemm3.9_1}
For $\xi_1 , \xi_2 \in Q^{\vee}$, 
$T^J_{\xi_1} \succeq T^J_{\xi_2}$ 
in 
$(W^J)_{\af}$
if and only if 
$[\xi_1]_{J^c} \succeq [\xi_2]_{J^c}$ 
in 
$Q^{\vee}$.
\item\label{Lemm3.9_2}
For 
$x , y \in W_{\af}$
and 
$\xi \in Q^{\vee}$, 
$x \succeq y$ 
in 
$W_{\af}$
if and only if
$x t_{\xi} \succeq yt_{\xi}$
in 
$W_{\af}$. 
In particular, we have 
$wt_{\xi} \succeq t_{\xi}$ 
in 
$W_{\af}$ 
for all 
$w \in W$.
\item\label{Lemm3.9_3}
For any 
$y \in W_{\af}$
and 
$\xi \in Q^{\vee}$, 
there exists 
$\vartheta \in Q^{\vee}$
such that 
$\vartheta \succeq \xi$ 
in 
$Q^{\vee}$
and 
$t_{\vartheta} \succeq y$
in 
$W_{\af}$.
\item\label{Lemm3.9_4}
Let 
$J,K \subset I$ 
be such that 
$J^c \cap K^c = \emptyset$. 
Then, 
for any 
$x , y \in W_{\af}$, 
there exist
$\vartheta_1 \in Q_J^{\vee}$ 
and 
$\vartheta_2 \in Q_K^{\vee}$
such that 
$\Pi^J(x t_{\vartheta_1}) = \Pi^{J}(x)$, 
$\Pi^K(y t_{\vartheta_2}) = \Pi^K(y)$ and 
$x t_{\vartheta_1} \succeq y t_{\vartheta_2}$
in
$W_{\af}$.
\item\label{Lemm3.9_5}
Let $J \subset I$ and $x , y \in W_{\af}$. 
If 
$x \succeq y$ in $W_{\af}$, 
then 
$\Pi^J(x) \succeq \Pi^J(y)$ in $(W^J)_{\af}$.
\end{enumerate}
\end{lemm}

\begin{proof}
\ref{Lemm3.9_1}: 
This is a special case of~\cite[Lemma~6.2.1]{INS16}.

\ref{Lemm3.9_2}: 
The assertion follows immediately from the formula 
$\ell^{\si}(x t_{\xi}) = \ell^{\si}(x) + 2(\xi,\rho)$. 

\ref{Lemm3.9_3}: 
Let 
$y = v t_{\zeta}$, 
$v \in W$ 
and 
$\zeta \in Q^{\vee}$, 
and let
$v = r_{i_1} r_{i_2} \cdots r_{i_l}$, 
$i_1 , i_2 , \ldots , i_l \in I_{\af}$, 
be a reduced expression. 
If we set 
$w_k = r_{i_1} r_{i_2} \cdots r_{i_k}$
and 
$\gamma_k = w_k \alpha_{i_k}$
for 
$k = 1,2,\ldots ,l$, then
\begin{align}
y = w_l t_{\zeta} 
\xrightarrow{\delta + \gamma_l} 
w_{l-1} t_{\zeta + \alpha_{i_l}^{\vee}} 
\xrightarrow{\delta + \gamma_{l-1}}
\cdots 
\xrightarrow{\delta + \gamma_1}
t_{\zeta + \alpha_{i_1}^{\vee} + \alpha_{i_2}^{\vee} + \cdots + \alpha_{i_l}^{\vee}}
\ \ 
\text{in}
\ 
\SiB ^{\emptyset}, 
\end{align}
which proves that 
$\vartheta' = \zeta + \alpha_{i_1}^{\vee} + \alpha_{i_2}^{\vee} + \cdots + \alpha_{i_l}^{\vee} \in Q^{\vee}$ 
satisfies 
$t_{\vartheta'} \succeq y$. 
We see that 
$\vartheta = \sum_{i \in I} \max \{ c_i(\xi) , c_i(\vartheta') \} \alpha_i^{\vee}$ 
satisfies 
$\vartheta \succeq \xi$, 
$\vartheta \succeq \vartheta'$, 
and hence 
$t_{\vartheta} \succeq t_{\vartheta'}$
by~\ref{Lemm3.9_1}.

\ref{Lemm3.9_4}: 
Assume that 
$x = wt_{\xi}$
with 
$w \in W$
and 
$\xi \in Q^{\vee}$. 
By~\ref{Lemm3.9_3}, 
there exists 
$\vartheta \in Q^{\vee}$ 
such that 
$\vartheta \succeq \xi$
and 
$t_{\vartheta} \succeq y$. 
If we set 
$\vartheta_1 = [\vartheta - \xi]_J$
and 
$\vartheta_2 =[\xi - \vartheta]_K$, 
then
$\Pi^J(x t_{\vartheta_1}) = \Pi^J(x)$
and 
$\Pi^K(y t_{\vartheta_2}) = \Pi^K(y)$. 
Moreover, we have 
$\xi + \vartheta_1 \succeq \vartheta + \vartheta_2$, 
because 
$\vartheta \succeq \xi$, 
$\xi + \vartheta_1 = [\xi]_{J^c} + [\vartheta]_{K^c} + [\vartheta]_{J \cap K}$
since 
$J = K^c \sqcup (J \cap K)$, and 
$\vartheta + \vartheta_2 = [\xi]_{J^c} + [\vartheta]_{K^c} + [\xi]_{J \cap K}$
since
$K = J^c \sqcup (J \cap K)$. 
Then~\ref{Lemm3.9_1}--\ref{Lemm3.9_2} shows that
\begin{align}
xt_{\vartheta_1} 
=
wt_{\xi + \vartheta_1}
\succeq
t_{\xi + \vartheta_1}
\succeq
t_{\vartheta + \vartheta_2}
=
t_{\vartheta} t_{\vartheta_2}
\succeq
yt_{\vartheta_2}.
\end{align}

\ref{Lemm3.9_5}: 
By induction on $\ell^{\si}(x) - \ell^{\si}(y)$, 
the assertion follows from~\cite[Lemma~6.1.1 for $K = \emptyset$]{INS16}.
\end{proof}

\begin{proof}[Proof of Theorem~\ref{thm:SMT}]
We first prove that~\ref{Theo3.4_2} and~\ref{Theo3.4_3} are equivalent. 
Clearly,~\ref{Theo3.4_2}~implies~\ref{Theo3.4_3}. 
We prove that~\ref{Theo3.4_3}~implies~\ref{Theo3.4_2}. 
The proof is by induction on 
$\# J_{\lambda}^c$. 
If 
$\# J_{\lambda}^c = 1$, 
then~\ref{Theo3.4_2} and~\ref{Theo3.4_3} are equivalent. 
Assume that 
$\# J_{\lambda}^c > 1$, 
$b \in \Breve{\mathcal{B}}(\lambda)$
satisfies~\ref{Theo3.4_3}, and 
$\Psi_{\lambda}(b) = \Psi^{(i)}_{\lambda}(b) \otimes \bigotimes_{j \in J_{\lambda}^c \setminus \{ i \}} \Psi^{(j)}_{\lambda}(b)$.
By~\ref{Theo3.4_3}, there exists a defining chain 
$(x_1 , \ldots , x_N)$
for 
$\Psi_{\lambda}^{(i)}(b)$. 
By induction hypothesis, 
there exists a defining chain 
$(y_1 , \ldots , y_M)$
for 
$\bigotimes_{j \in J_{\lambda}^c \setminus \{ i \}} \Psi_{\lambda}^{(j)}(b)$. 
Applying Lemma~\ref{lem:easy}~\ref{Lemm3.9_4} to
$x = x_N$, $y = y_1$, $J = I \setminus \{ i \}$, and $K = J_{\lambda} \cup \{ i \}$
to obtain 
$\vartheta_1 \in Q_{I \setminus \{ i \}}^{\vee}$
and 
$\vartheta_2 \in Q_{J_{\lambda} \cup \{ i \}}^{\vee}$
such that 
$x_N t_{\vartheta_1} \succeq y_1 t_{\vartheta_2}$. 
By Lemma~\ref{lem:easy}~\ref{Lemm3.9_2}, 
we conclude that 
$(x_1 t_{\vartheta_1} , \ldots , x_N t_{\vartheta_1} , y_1 t_{\vartheta_2} , \ldots , y_M t_{\vartheta_2})$
is a defining chain for 
$\Psi_{\lambda}(b)$.

We next prove that~\ref{Theo3.4_1} and~\ref{Theo3.4_2} are equivalent. 
The proof is completed by showing that 
\begin{enumerate}[label=(\roman*)]
\item\label{proof_theo3.4_i}
$\Breve{\mathbb{S}}^{\si}(\lambda)$ 
is stable under the Kashiwara operators, and 
\item\label{proof_theo3.4_ii}
each connected component of 
$\Breve{\mathbb{S}}^{\si}(\lambda)$ 
contains 
$\Psi_{\lambda}(z^{-\bmrho{\rho}} \Tilde{u}_{\lambda})$
for some 
$\bmrho{\rho} \in \Par (\lambda)$, 
\end{enumerate}
because 
$\Psi_{\lambda} \circ \Phi_{\lambda \mid q=0}^{\LT } : \mathcal{B}(\lambda) \rightarrow \Breve{\mathbb{B}}^{\si}(\lambda)$ 
is a strict embedding of $\U$-crystals, and 
$\Psi_{\lambda}(z^{-\bmrho{\rho}} \Tilde{u}_{\lambda}) \in \Breve{\mathbb{S}}^{\si}(\lambda)$
for every 
$\bmrho{\rho} \in \Par (\lambda)$; 
indeed, 
by Lemmas~\ref{lem:SiLS_W-action}~\ref{Lemm3.8_1} and~\ref{lem:easy}~\ref{Lemm3.9_1}, 
we have a defining chain
$\left( t_{\rho^{(i)}_1 \alpha_i^{\vee}} , \ldots , t_{\rho^{(i)}_{m_i-1} \alpha_i^{\vee}} , e \right)$ 
for
$\Psi_{\lambda}^{(i)}(z^{-\bmrho{\rho}} \Tilde{u}_{\lambda})$
for every 
$\bmrho{\rho} = (\rho^{(i)})_{i \in I} \in \Par (\lambda)$, 
with 
$\rho^{(i)} = (\rho^{(i)}_1 \geq \cdots \geq \rho^{(i)}_{m_i-1} \geq 0)$, 
$i \in I$, 
and hence there exists a defining chain for 
$\Psi_{\lambda}(z^{-\bmrho{\rho}} \Tilde{u}_{\lambda})$
by the implication~\ref{Theo3.4_3}~$\Rightarrow$~\ref{Theo3.4_2}.

We prove~\ref{proof_theo3.4_i} only for the action of $e_j$, $j \in I_{\af}$; 
the proof for $f_j$ is similar. 
Let 
$\eta = \bigotimes_{i \in J_{\lambda}^c} \eta^{(i)}\in \Breve{\mathbb{S}}^{\si}(\lambda)$, 
with 
$\eta^{(i)}\in \Breve{\mathbb{S}}^{\si}(m_i \varpi_i)$, 
$i \in J_{\lambda}^c$. 
By tensor product rule 
and the implication~\ref{Theo3.4_3}~$\Rightarrow $~\ref{Theo3.4_2}, 
we only need to show that, for each 
$i \in J_{\lambda}^c$, 
if 
$e_j \eta^{(i)} \neq \bm{0}$, 
then 
$e_j \eta^{(i)} \in \Breve{\mathbb{S}}^{\si}(m_i \varpi_i)$. 
Write 
$\eta^{(i)} = \bigotimes_{\nu = 1}^{m_i} \eta^{(i)}_{\nu}$, 
with 
$\eta^{(i)}_{\nu} \in \mathbb{B}^{\si}(\varpi_i)$, 
$1 \leq \nu \leq m_i$, 
and let
$(x^{[\nu]}_s)_{1 \leq s \leq l_{\nu}, 1 \leq \nu \leq m_i}$
be a defining chain for 
$\eta^{(i)}$; 
by Lemma~\ref{lem:easy}~\ref{Lemm3.9_5}, 
we may assume that 
$x^{[\nu]}_s \in (W^{I \setminus \{ i \}})_{\af}$, 
$1 \leq s \leq l_{\nu}$, 
$1 \leq \nu \leq m_i$.
By tensor product rule, 
$e_j \eta^{(i)} = 
\eta^{(i)}_1 \otimes \cdots \otimes e_j \eta^{(i)}_{\nu} \otimes \cdots \otimes \eta^{(i)}_{m_i}$
for some
$1 \leq \nu \leq m_i$. 
Let 
$1 \leq p < q \leq l_{\nu}$
be as in~\eqref{eq:e_i} for $\eta^{(i)}_{\nu}$. 
It follows from 
Proposition~\ref{prop:conc=tensor} 
and~\cite[Lemma~4.1.6]{INS16} 
that the tuple
\begin{align}
\left( 
x^{[1]}_1 , \ldots , x^{[\nu-1]}_{l_{\nu -1}}, 
x^{[\nu]}_1 , \ldots , x^{[\nu]}_p , r_i x^{[\nu]}_p , \ldots , 
r_i x^{[\nu]}_q , x^{[\nu]}_{q+1} , \ldots , x^{[\nu]}_{l_{\nu}} , 
x^{[\nu +1]}_1 , \ldots , x^{[m_i]}_{l_{m_i}}
\right)
\end{align}
is a defining chain for 
$e_j \eta^{(i)}$, 
and hence 
$e_j \eta^{(i)} \in \Breve{\mathbb{S}}^{\si}(m_i \varpi_i)$. 

Finally, we prove~\ref{proof_theo3.4_ii}. 
Let 
$C$ 
be an arbitrary connected component of 
$\Breve{\mathbb{S}}^{\si}(\lambda)$; 
we see from~\ref{proof_theo3.4_i} that 
$C$ 
is a connected component of 
$\Breve{\mathbb{B}}^{\si}(\lambda)$. 
By Lemma~\ref{lem:SiLS_W-action}~\ref{Lemm3.8_2}, 
$C$ 
contains an element of the form 
$\eta = \bigotimes_{i \in J_{\lambda}^c} \bigotimes_{\nu = 1}^{m_i} \left(T^{I \setminus \{ i \}}_{\rho^{(i)}_{\nu} \alpha_i^{\vee}};0,1 \right)$, 
with 
$\rho^{(i)}_{m_i} = 0$, 
$i \in J_{\lambda}^c$. 
Since there exists a defining chain for $\eta$, 
we see from Lemma~\ref{lem:easy}~\ref{Lemm3.9_1} that 
$\rho^{(i)} = (\rho^{(i)}_1 \geq \rho^{(i)}_2 \geq \cdots \geq \rho^{(i)}_{m_i -1})$
is a partition of length less than $m_i$ 
for each $i \in I$; 
here, we set 
$\rho^{(i)} = \emptyset$ 
if 
$i \in J_{\lambda}$. 
Hence 
$\bmrho{\rho} = (\rho^{(i)})_{i \in I} \in \Par (\lambda)$. 
We have 
$\Psi_{\lambda}(z^{-\bmrho{\rho}} \Tilde{u}_{\lambda}) = \eta \in C$
by Lemma~\ref{lem:SiLS_W-action}~\ref{Lemm3.8_1}, 
which proves~\ref{proof_theo3.4_ii}. 
\end{proof}

\section{Semi-infinite Young tableaux} \label{Section:SiYT}

Throughout this section, 
we will make the following assumptions: 
$\Fg_{\af}$ 
is of type 
$A_{n-1}^{(1)}$
(see~\cite[\S~4.8 TABLE Aff1]{Kac}), 
and 
$I = \{ 1, 2, \ldots ,n-1 \}$ 
satisfies 
\begin{align}
(\alpha_i , \alpha_j)
=
\begin{cases}
2 & \text{if} \ i=j, \\
-1 & \text{if} \ i-j \equiv \pm 1 \mod n, \\
0 & \text{otherwise}
\end{cases}
\end{align}
for all $i , j \in I_{\af} = \{ 0 \} \sqcup I$. 
In this case, 
$\alpha_i^{\vee} = \alpha_i$, $i \in I_{\af}$, 
and hence 
$Q^{\vee} = Q$.
We sometimes think of $W$ as the permutation group of 
$\{ 1, 2, \ldots ,n \}$, namely the symmetric group of degree $n$, 
where $r_i$, $i \in I$, acts as the transposition $(i \ i+1)$. 
Observe that this action extends to the $W_{\af}$-action, 
where 
$r_0$ 
acts as the transposition $(n \ 1)$; 
note that each 
$t_{\xi}$, $\xi \in Q$, 
acts as the identity. 

\subsection{Semi-infinite Young tableaux and isomorphism theorem} \label{subsection:SiYT}

We identify each element $\lambda = \sum_{i \in I} m_i \varpi_i \in P^+$ 
with the Young diagram such that the number of the columns of length $i$ is $m_i$ for each $i \in I$. 
A column-strict tableau of shape 
$\lambda \in P^+$ 
with entries in 
$\{ 1,2,\ldots ,n \}$
is, by definition, an assignment of a number in 
$\{ 1,2,\ldots ,n \}$
to each box of the Young diagram 
$\lambda$
such that the entries are strictly increasing 
from top to bottom in each column.
Let $\CST (\lambda)$ 
be the set of column-strict tableaux of shape $\lambda$ 
with entries in $\{ 1,2, \ldots ,n \}$. 
For a tuple 
$(\mathsf{T}_1 , \mathsf{T}_2 , \ldots , \mathsf{T}_M)$
of column-strict tableaux of one-column shapes, let 
$\prod_{\nu = 1}^M \mathsf{T}_{\nu}
= \mathsf{T}_1 \mathsf{T}_2 \cdots \mathsf{T}_M$
denote the column-strict tableau whose $\nu$-th column is 
$\mathsf{T}_{\nu}$. 
For $\mathsf{T} \in \CST (\varpi_i)$, 
let 
$\mathsf{T}(s) \in \{ 1,2,\ldots , n \}$, $1 \leq s \leq i$, 
denote the $s$-th entry (from top) of 
$\mathsf{T}$. 

\begin{rema} \label{rem:Young}
In this paper, 
we consider a Young diagram 
as a collection of boxes, 
arranged in right-justified rows, 
with a weakly decreasing number of boxes in each row 
from top to bottom.
For example, the Young diagram 
$\lambda = 5\varpi_1 + 3\varpi_2 + 4 \varpi_3 + 2 \varpi_4 + \varpi_6$
is as follows:
\begin{align}
\ytableaushort{
\ \ \ \ \ \ \ \ \ \ \ \ \ \ \ , 
\none \none \none \none \none \ \ \ \ \ \ \ \ \ \ ,
\none \none \none \none \none \none \none \none \ \ \ \ \ \ \ , 
\none \none \none \none \none \none \none \none \none \none \none \none \ \ \ ,
\none \none \none \none \none \none \none \none \none \none \none \none \none \none \ ,
\none \none \none \none \none \none \none \none \none \none \none \none \none \none \ }
\end{align}
\end{rema}

\begin{defi} \label{def:SiYT}\ \\[-1.2em]
\begin{enumerate}[label=(\arabic*)]
\item\label{Defi4.2_1}
Define the partial order 
$\preceq$
on 
$\CST (\varpi_i) \times \BZ$
as follows: 
for 
$(\mathsf{T},c) , (\mathsf{T}' , c') \in \CST (\varpi_i) \times \BZ$, 
set 
$(\mathsf{T},c) \preceq (\mathsf{T}' , c')$
if 
\begin{align}
\bigl(
c \leq c'
\bigr)
\ \ \text{and} \ \ 
\bigl(
\mathsf{T}(u) \leq \mathsf{T}'(u + c' - c)
\ \text{if} \ 
1 \leq u \leq i- c' + c
\bigr).
\end{align}
\item\label{Defi4.2_2}
Let 
$\lambda = \sum_{i \in I} m_i \varpi_i \in P^+$, 
$m_i \in \BZ_{\geq 0}$, $i \in I$, 
and 
$N = \sum_{i \in I} m_i$. 
Let 
\begin{align}
\mathbb{T} = 
\left(
\prod_{i=1}^{n-1} \prod_{\nu = 1}^{m_i}
\mathsf{T}_{\nu}^{(i)} , 
(c_{\nu}^{(i)})_{1 \leq \nu \leq m_i , 1 \leq i \leq n-1}
\right) 
\in \CST (\lambda) \times \BZ^N ,
\end{align}
where 
$\mathsf{T}^{(i)}_{\nu} \in \CST (\varpi_i)$ 
and 
$c^{(i)}_{\nu} \in \BZ$
for 
$1 \leq \nu \leq m_i$, $i \in I$. 
We call $\mathbb{T}$ a semi-infinite Young tableau of shape $\lambda$
if 
\begin{align}
(\mathsf{T}^{(i)}_1 , c_1^{(i)}) \succeq (\mathsf{T}^{(i)}_2 , c_2^{(i)}) \succeq \cdots \succeq (\mathsf{T}^{(i)}_{m_i} , c_{m_i}^{(i)})
\ \text{in} \ 
\CST (\varpi_i) \times \BZ
\ \text{for every} \ 
i \in I.
\end{align}
Let 
$\mathbb{Y}^{\si}(\lambda)$ 
be the set of semi-infinite Young tableaux of shape $\lambda$; 
note that 
$\mathbb{Y}^{\si}(\varpi_i) = \CST (\varpi_i) \times \BZ$.
\end{enumerate}
\end{defi}

In \S~\ref{Subsection:Y(pi_i)}, 
we define a $\U$-crystal structure on 
$\mathbb{Y}^{\si}(\varpi_i)$, 
and prove the next proposition. 

\begin{prop} \label{prop:aff-varpi_i}
Let $i \in I$.
\begin{enumerate}[label=(\arabic*)]
\item\label{Prop4.3_1}
There exists a unique isomorphism
$\Upsilon_i : \mathcal{B}(\varpi_i) \rightarrow \mathbb{Y}^{\si}(\varpi_i)$
of $\U$-crystals. 
\item\label{Prop4.3_2}
We have 
$\mathcal{B}(\varpi_i) = \{ u_x := S_x u_{\varpi_i} \mid x \in (W^{I \setminus \{ i \}})_{\af} \}$, 
and the map 
$(W^{I \setminus \{ i \}})_{\af} \rightarrow \mathcal{B}(\varpi_i)$, $x \mapsto u_x$, 
is bijective. In particular, 
$V(\varpi_i)$ 
is a minuscule representation of $\U$. 
\end{enumerate}
\end{prop}

\begin{rema} \label{rem:aff-varpi_i}
It follows from Theorem~\ref{thm:SiLS-isom}~\ref{Theo2.7_2} and Proposition~\ref{prop:aff-varpi_i} that 
$\mathbb{B}^{\si}(\varpi_i) = \{ (x ; 0,1) \mid x \in (W^{I \setminus \{ i \}})_{\af} \}$, 
and the map 
$\mathcal{B}(\varpi_i) \rightarrow \mathbb{B}^{\si}(\varpi_i)$, 
$u_x \mapsto (x;0,1)$, 
$x \in (W^{I \setminus \{ i \}})_{\af}$, 
equals the isomorphism in Theorem~\ref{thm:SiLS-isom}~\ref{Theo2.7_2}.
\end{rema}

Let 
$\lambda = \sum_{i \in I} m_i \varpi_i \in P^+$
and 
$N = \sum_{i \in I} m_i$. 
We have a bijection from 
$\bigotimes_{i=1}^{n-1} \mathbb{Y}^{\si}(\varpi_i)^{\otimes m_i}$
to
$\CST (\lambda) \times \BZ^N$
defined by
\begin{align} \label{eq:Y->CST}
\bigotimes_{i=1}^{n-1} \bigotimes_{\nu = 1}^{m_i} 
\left( \mathsf{T}_{\nu}^{(i)} , c_{\nu}^{(i)} \right)
\longmapsto 
\left(
\prod_{i=1}^{n-1} \prod_{\nu = 1}^{m_i}
\mathsf{T}_{\nu}^{(i)} , 
(c_{\nu}^{(i)})_{1 \leq \nu \leq m_i , 1 \leq i \leq n-1}
\right) ,
\end{align}
where 
$(\mathsf{T}_{\nu}^{(i)} , c_{\nu}^{(i)}) \in \mathbb{Y}^{\si}(\varpi_i)$, 
$1 \leq \nu \leq m_i$, $i \in I$. 
Define a $\U$-crystal structure on 
$\CST (\lambda) \times \BZ^N$ 
to be such that the map~\eqref{eq:Y->CST} is an isomorphism of $\U$-crystals. 
From now on we assume that 
$\Breve{\mathcal{B}}(\lambda) 
=
\bigotimes_{i=1}^{n-1} \mathcal{B}(\varpi_i)^{\otimes m_i}
=
\mathcal{B}(\varpi_1)^{\otimes m_1} \otimes \cdots \otimes \mathcal{B}(\varpi_{n-1})^{\otimes m_{n-1}}$.

\begin{theo} \label{thm:SiYT}
Let 
$\lambda = \sum_{i \in I} m_i \varpi_i \in P^+$
and 
$N = \sum_{i \in I} m_i$. 
Then, 
$\mathbb{Y}^{\si}(\lambda)$ 
equals the image of the composition of the maps
\begin{align} \label{eq:emb-SiYT}
\mathcal{B}(\lambda) 
\xrightarrow[\S~\ref{Subsection:LT}]{\Phi_{\lambda \mid q=0}^{\LT }}
\Breve{\mathcal{B}}(\lambda) 
= 
\bigotimes_{i=1}^{n-1} \mathcal{B}(\varpi_i)^{\otimes m_i}
\xrightarrow[\begin{subarray}{c}\ \textup{Proposition} \\ %\hspace{3.5mm} 
{\ref{prop:aff-varpi_i}}~\ref{Prop4.3_1} \end{subarray}]{\bigotimes_{i=1}^{n-1} \Upsilon_i^{\otimes m_i}} 
\bigotimes_{i=1}^{n-1} \mathbb{Y}^{\si}(\varpi_i)^{\otimes m_i}
\xrightarrow[\eqref{eq:Y->CST}]{\cong}
\CST (\lambda) \times \BZ^N .
\end{align}
\end{theo}

Since the map~\eqref{eq:emb-SiYT} is a strict embedding of $\U$-crystals, 
we have the following.

\begin{coro} \label{cor:isom}
Let 
$\lambda = \sum_{i \in I} m_i \varpi_i \in P^+$
and 
$N = \sum_{i \in I} m_i$. 
Then, 
$\mathbb{Y}^{\si}(\lambda)$ 
is stable under the Kashiwara operators 
on 
$\bigotimes_{i=1}^{n-1} \mathbb{Y}^{\si}(\varpi_i)^{\otimes m_i} \cong \CST (\lambda) \times \BZ^N$, 
and is isomorphic, as a $\U$-crystal, to the crystal basis 
$\mathcal{B}(\lambda)$.
\end{coro}

Theorem~\ref{thm:SiYT} follows from 
Theorem~\ref{thm:SMT}, Definition~\ref{def:SiYT}, Remark~\ref{rem:aff-varpi_i}, 
and the following tableau criterion for the semi-infinite Bruhat order. 

\begin{theo} \label{thm:tab-cri}
Let 
$i \in I$ 
and 
$x,y \in (W^{I \setminus \{ i \}})_{\af}$. 
Write 
$\Upsilon_i(u_x) = (\mathsf{T},c)$
and 
$\Upsilon_i(u_y) = (\mathsf{T}',c')$. 
The following conditions are equivalent:
\begin{enumerate}[label=(\arabic*)]
\item\label{Theo4.7_1}
$x \preceq y$ 
in 
$(W^{I \setminus \{ i \}})_{\af}$. 
\item\label{Theo4.7_2}
$c \leq c'$ 
and 
$\mathsf{T}(u) \leq \mathsf{T}'(u + c' - c)$ 
if 
$1 \leq u \leq i - c' + c$. 
\end{enumerate}
\end{theo}

Theorem~\ref{thm:tab-cri} will be proved in \S~\ref{Subsection:tab-cri}.

\subsection{Explicit description of \texorpdfstring{$(W^J)_{\mathrm{af}}$}{(W J) af}}

In this subsection, 
following~\cite[\S~3]{LNSSS15}, 
we give an explicit description of 
$(W^J)_{\af}$ 
for later use. 

We take and fix $J = \bigsqcup_{m=1}^k I_m \subset I$, 
where $I_1 , I_2 , \ldots , I_k$ are the sets of vertices 
of the connected components of the Dynkin diagram of $\Delta_J$; 
note that 
$\Delta_J = \bigsqcup_{m=1}^k \Delta_{I_m}$ 
and each 
$\Delta_{I_m}$, $1 \leq m \leq k$, 
is of finite type $A$. 
Set 
$(I_m)_{\af} = \{ 0 \} \sqcup I_m \subset I_{\af}$, 
$1 \leq m \leq k$.
For $1 \leq s \leq t \leq n-1$, set 
$\alpha_{s,t} = \sum_{i=s}^t \alpha_i$; 
note that
$\alpha_s = \alpha_{s,s}$.
It follows that 
\begin{align} \label{eq:root-A}
\Delta = \{ \pm \alpha_{s,t} \mid 1 \leq s \leq t \leq n-1 \} .
\end{align}
Set 
\begin{align}
Q^J = \{ \xi \in Q \mid (\xi,\alpha) \in \{ -1,0\} \ \text{for all} \ \alpha \in \Delta_J^+ \} .
\end{align}

\begin{lemm}[{\cite[Equation~(3.6)]{LNSSS15}}] \label{lem:J-ad}
For each $\xi \in Q$ there exist a unique
$\phi_J(\xi) \in Q_J$ 
and a unique 
$(j_1, j_2 , \ldots , j_k) \in \prod_{m=1}^k (I_m)_{\af}$ 
such that 
\begin{align} \label{eq:J-ad}
\xi + \phi_J(\xi) + \sum_{m=1}^k \varpi_{j_m} \in \bigoplus_{i \in I \setminus J} \BZ \varpi_i \oplus \BC \delta.
\end{align}
In particular, $\xi + \phi_J(\xi) \in Q^J$ for any $\xi \in Q$, 
and hence $Q^J$ is a complete system of coset representatives for $Q/Q_J$. 
\end{lemm}

For a subset 
$K \subset I$, 
let 
$w_0^K$ 
be the longest element of 
$W_K$.
For $j_m \in (I_m)_{\af}$, set 
\begin{align}
v_{j_m}^{I_m} = w_0^{I_m} w_0^{I_m \setminus \{ j_m \}} \in W_{I_m} \subset W_J; 
\end{align}
note that $v^{I_m}_0 = e$.
For $\xi \in Q$, define
\begin{align} \label{eq:z}
z_{\xi} = z_{\xi}^J = v_{j_1}^{I_1} v_{j_2}^{I_2} \cdots v_{j_k}^{I_k} \in W_J,
\end{align}
where 
$(j_1, j_2 , \ldots , j_k) \in \prod_{m=1}^k (I_m)_{\af}$, 
satisfying~\eqref{eq:J-ad} for $\xi$, 
is determined uniquely by Lemma~\ref{lem:J-ad}; 
note that 
$z_{\xi} = z_{\zeta}$ 
if 
$\xi \equiv \zeta \mod Q_J$.

\begin{lemm}[{\cite[Lemma~3.7]{LNSSS15}}] \label{lem:(W^J)_af}
We have 
$T_{\xi} = \Pi^J(t_{\xi}) = z_{\xi} t_{\xi + \phi_J(\xi)}$
for every $\xi \in Q$. 
Therefore, by Lemma~\ref{lem:Pi}, 
$\Pi^J(wt_{\xi}) = \lfloor w \rfloor z_{\xi} t_{\xi + \phi_J(\xi)}$
for every $w \in W$ and $\xi \in Q$, and we have a bijection
$W^J \times Q^J \rightarrow (W^J)_{\af}$, 
$(w,\xi) \mapsto wT_{\xi}$. 
In particular, 
\begin{align}
(W^J)_{\af} 
= 
\{ wT_{\xi} = w z_{\xi}t_{\xi} \mid w \in W^J , \ \xi \in Q^J \} .
\end{align}
\end{lemm}

\subsection{Crystal structure on \texorpdfstring{$\mathbb{Y}^{\si}(\varpi_i)$}{Y (infty/2)(varpi i)}} \label{Subsection:Y(pi_i)}

In this subsection, 
we define a $\U$-crystal structure on 
$\mathbb{Y}^{\si}(\varpi_i)$, 
and give a proof of Proposition~\ref{prop:aff-varpi_i}. 

We know from~\cite[Lemma~2.4.7]{BB}
that
\begin{align} \label{eq:Gr}
W^{I \setminus \{ i \}} 
\Mk = \Mk 
\{ w \in W \mid w(1) \Mk <\Mk w(2) \Mk < \Mk \cdots \Mk <\Mk w(i), \text{ and } w(i+1) \Mk <\Mk w(i+2) 
\Mk <\Mk \cdots \Mk < \Mk w(n) \} .
\end{align}
For 
$w \in W^{I \setminus \{ i \}}$, 
set
\ytableausetup{mathmode,boxsize=9mm}
\begin{align}
\mathsf{T}_w = 
\begin{ytableau}
w(1) \\ w(2) \\ \vdots \\ w(i)
\end{ytableau}
\ \in 
\CST (\varpi_i) .
\end{align}
By~\eqref{eq:Gr}, 
we have 
$\CST (\varpi_i) = \{ \mathsf{T}_w \mid w \in W^{I \setminus \{ i \}} \}$ 
and the map
$W^{I \setminus \{ i \}} \rightarrow \CST (\varpi_i)$, $w \mapsto \mathsf{T}_w$, 
is bijective. 
Let 
$c_i(\xi)$ 
be the coefficient of $\alpha_i \left(= \alpha_i^{\vee}\right)$ in $\xi \in Q$. 
It follows from Lemma~\ref{lem:J-ad} that 
$Q^{I \setminus \{ i \}} = \{ c\alpha_i + \phi_{I\setminus \{ i \}}(c\alpha_i) \mid c \in \BZ \}$, 
and the maps 
$\BZ \rightarrow Q^{I \setminus \{ i \}}$, $c \mapsto c\alpha_i + \phi_{I\setminus \{ i \}}(c\alpha_i)$, 
and
$Q^{I \setminus \{ i \}} \rightarrow \BZ$, $\xi \mapsto c_i(\xi)$, 
are inverses of each other. 
We have thus proved that the map
\begin{align} \label{eq:W->Y}
\mathcal{Y}_i : 
(W^{I \setminus \{ i \}})_{\af} \rightarrow \mathbb{Y}^{\si}(\varpi_i),\ 
wT_{\xi} \mapsto 
\mathcal{Y}_i(wT_{\xi}) = 
\left( \mathsf{T}_w , c_i(\xi) \right),
\end{align}
is bijective, where 
$w \in W^{I \setminus \{ i \}}$
and 
$\xi \in Q^{I \setminus \{ i \}}$
(see Lemma~\ref{lem:(W^J)_af}).

Following~\cite[\S~3.7]{S} (see also~\cite[\S~4.1]{KMN}), we equip the set 
$\CST (\varpi_i)$
with a $\U'$-crystal structure as follows: 
let 
$\mathsf{T} \in \CST (\varpi_i)$. 
For 
$k \in \{ 1,2,\ldots ,n \}$, 
write 
$k \in \mathsf{T}$ 
if 
$\mathsf{T}(s) = k$
for some 
$1 \leq s \leq i$.
\begin{enumerate}[label=(\roman*)]
\item
Define 
$\wt (\mathsf{T}_w) = w\varpi_i$ (mod $\BC \delta$) 
for 
$w \in W^{I \setminus \{ i \}}$. 
\item
For $j \in I$, 
if 
$\mathsf{T}(s) = j+1$ 
and 
$j \notin \mathsf{T}$, 
then we define 
$e_j \mathsf{T} \in \CST (\varpi_i)$ 
to be such that 
$(e_j \mathsf{T})(s) = j$ 
and 
$(e_j \mathsf{T})(u) = \mathsf{T}(u)$ 
for 
$1 \leq u \leq i$, $u \neq s$. 
\item
If 
$1 \in \mathsf{T}$
and 
$n \notin \mathsf{T}$, 
then we define 
$e_0 \mathsf{T} \in \CST (\varpi_i)$ 
to be such that 
$(e_0 \mathsf{T})(i) = n$ 
and 
$(e_0 \mathsf{T})(u) = \mathsf{T}(u+1)$ 
for 
$1 \leq u \leq i-1$.
\item
Otherwise, we set 
$e_j \mathsf{T} = \bm{0}$
for $j \in I_{\af}$.
\item
For $j \in I$, 
if
$\mathsf{T}(s) = j$ 
and 
$j+1 \notin \mathsf{T}$, 
then we define 
$f_j \mathsf{T} \in \CST (\varpi_i)$ 
to be such that 
$(f_j \mathsf{T})(s) = j+1$ 
and 
$(f_j \mathsf{T})(u) = \mathsf{T}(u)$ 
for 
$1 \leq u \leq i$, $u \neq s$. 
\item
If 
$1 \notin \mathsf{T}$ 
and 
$n \in \mathsf{T}$, 
then we define 
$f_0 \mathsf{T} \in \CST (\varpi_i)$ 
to be such that 
$(f_0 \mathsf{T})(1) = 1$ 
and 
$(f_0 \mathsf{T})(u) = \mathsf{T}(u-1)$ 
for 
$2 \leq u \leq i$.
\item
Otherwise, we set 
$f_j \mathsf{T} = \bm{0}$
for $j \in I_{\af}$. 
\item
Define
\begin{align} \label{eq:CST}
\varepsilon_j(\mathsf{T}) 
=
\begin{cases}
1 & \text{if} \ e_j \mathsf{T} \neq \bm{0}, \\%[1mm]
0 & \text{if} \ e_j \mathsf{T} = \bm{0},
\end{cases}
&&
\varphi_j(\mathsf{T}) 
=
\begin{cases}
1 & \text{if} \ f_j \mathsf{T} \neq \bm{0}, \\%[1mm]
0 & \text{if} \ f_j \mathsf{T} = \bm{0}.
\end{cases}
\end{align}
\end{enumerate}

\begin{rema} \label{rem:KR}
The 
$\U'$-crystal 
$\CST (\varpi_i)$ 
defined above is isomorphic to the crystal basis of the $\U'$-module 
$W(\varpi_i)$ (see \S~\ref{Subsection:extremal}).
Indeed, we see from~\cite[Theorem~5.17~(ix)]{K02}
(see also~\cite[Remark~3.3]{N}) 
that 
$W(\varpi_i)$ 
is isomorphic to a Kirillov--Reshetikhin module, 
whose crystal basis is a perfect crystal of level $1$ 
in the sense of~\cite[Definition~1.1.1]{KMN}. 
It follows that the $\U'$-crystal 
$\CST (\varpi_i)$ 
and the crystal basis of 
$W(\varpi_i)$ 
satisfy the conditions in~\cite[Proposition~1.2.1 for $l = 1$]{KMN}, 
and hence they must be isomorphic to each other.
\end{rema}

The set 
$\mathbb{Y}^{\si}(\varpi_i) = \CST (\varpi_i) \times \BZ$ 
can be identified with the affinization of the $\U'$-crystal 
$\CST (\varpi_i)$
in the sense of~\cite[\S~4.2]{K02}. 
We have thus obtained a $\U$-crystal structure on 
$\mathbb{Y}^{\si}(\varpi_i)$
as follows: 
for 
$w \in W^{I \setminus \{ i \}}$, 
$c \in \BZ$, 
$\mathsf{T} \in \CST (\varpi_i)$, 
and 
$j \in I_{\af}$, 
\begin{align} \label{eq:CST*Z}
\begin{cases}
\wt (\mathsf{T}_w,c) = w\varpi_i - c\delta , \\%[1mm]
e_j (\mathsf{T},c) = (e_j \mathsf{T} , c - \delta_{j,0}), \ 
f_j (\mathsf{T},c) = (f_j \mathsf{T} , c + \delta_{j,0}), \\%[1mm]
\varepsilon_j(\mathsf{T},c) = \varepsilon_j(\mathsf{T}), \ 
\varphi_j(\mathsf{T},c) = \varphi_j(\mathsf{T}) ;
\end{cases}
\end{align}
we understand that 
$(\bm{0},c) = \bm{0}$. 
By~\eqref{eq:lv0-act},~\eqref{eq:W->Y}, and~\eqref{eq:CST*Z}, 
we have 
$\wt (\mathcal{Y}_i(x)) = x\varpi_i$ 
for all 
$x \in (W^{I \setminus \{ i \}})_{\af}$.

\begin{proof}[Proof of Proposition~\ref{prop:aff-varpi_i}]
\ref{Prop4.3_1}: 
Since 
$\mathbb{Y}^{\si}(\varpi_i)$ 
is isomorphic to the affinization of the crystal basis of 
$W(\varpi_i)$ (see Remark~\ref{rem:KR}), 
we see from~\cite[Proposition~5.4~(ii) and Theorem~5.17~(vii)]{K02}
that 
$\mathcal{B}(\varpi_i)$
is isomorphic, as a $\U$-crystal, to 
$\mathbb{Y}^{\si}(\varpi_i)$. 
Note that 
$(\mathsf{T}_e , 0) \in \mathbb{Y}^{\si}(\varpi_i)$ 
(and hence $u_{\varpi_i} \in \mathcal{B}(\varpi_i)$)
is a unique element of weight $\varpi_i$; 
indeed, by~\eqref{eq:CST*Z}, 
$\wt (\mathsf{T}_w,c) = \varpi_i$ 
holds if and only if 
$w = e$ and $c = 0$.
This and the connectedness of 
$\mathcal{B}(\varpi_i)$ 
(see~\cite[Proposition~5.4~(ii)]{K02})
prove the uniqueness of the isomorphism between 
$\mathcal{B}(\varpi_i)$
and 
$\mathbb{Y}^{\si}(\varpi_i)$.

\ref{Prop4.3_2}: 
By~\ref{Prop4.3_1} and~\eqref{eq:CST}--\eqref{eq:CST*Z}, we see that 
$\varepsilon_j(b) , \varphi_j(b) \in \{ 0,1 \}$ 
for all 
$j \in I_{\af}$ 
and 
$b \in \mathcal{B}(\varpi_i)$. 
Hence 
$e_j b = S_{r_j} b$
(\resp $f_j b = S_{r_j} b$) 
if 
$e_j b \neq \bm{0}$
(\resp $f_j b\neq \bm{0}$)
for $j \in I_{\af}$ and $b \in \mathcal{B}(\varpi_i)$. 
Since 
$\mathcal{B}(\varpi_i)$
is connected 
(\cite[Proposition~5.4~(ii)]{K02}), 
this proves that the action of 
$W_{\af}$ 
is transitive. 
We have 
$\{ x \in W_{\af} \mid S_x u_{\varpi_i} = u_{\varpi_i} \} = (W_{I \setminus \{ i \}})_{\af}$
(see~\eqref{eq:W_J_af}) 
by~\cite[Proposition~5.1.1]{INS16}, and hence 
$\mathcal{B}(\varpi_i) = \{ S_x u_{\varpi_i} \mid x \in (W^{I \setminus \{ i \}})_{\af} \}$
by Lemma~\ref{lem:Waf-decomp}. 
\end{proof}

\subsection{Tableau criterion for semi-infinite Bruhat order} \label{Subsection:tab-cri}

This subsection is devoted to the proof of Theorem~\ref{thm:tab-cri}.

We take and fix 
$i \in I$. 
It is easily seen from~\eqref{eq:root-A} that
\begin{align} \label{eq:Delta^i}
\Delta^+ \setminus \Delta_{I \setminus \{ i \}}^+
=
\{ \alpha_{s,t} \mid 1 \leq s \leq i \leq t \leq n-1 \} .
\end{align}
We have 
\begin{equation}
\Upsilon_i(u_x) = \mathcal{Y}_i(x)\text{ for all }x \in (W^{I \setminus \{ i \}})_{\af}
\end{equation} 
(see Proposition~\ref{prop:aff-varpi_i}~\ref{Prop4.3_1} and~\eqref{eq:W->Y}) 
because both elements are of weight $x\varpi_i$, 
and there is only one element of weight $x\varpi_i$ in 
$\mathbb{Y}^{\si}(\varpi_i)$.

\begin{prop} \label{prop:tab-cri}
Let 
$w \in W^{I \setminus \{ i \}}$, 
$\xi \in Q^{I \setminus \{ i \}}$, 
$\beta = w\gamma + \chi \delta \in \Delta_{\af}^+$, 
$\gamma = \alpha_{s,t} \in \Delta^+ \setminus \Delta_{I \setminus \{ i \}}^+$, 
$1 \leq s \leq i \leq t \leq n-1$, and
$\chi \in \BZ_{\geq 0}$. 
Write 
$\mathcal{Y}_i (wT_{\xi}) = (\mathsf{T},c)$
and 
$\mathcal{Y}_i \left( \Pi^{I \setminus \{ i \}}(r_{\beta}wT_{\xi}) \right) = (\mathsf{T}',c')$. 
Then 
$r_{\beta}wT_{\xi} \in (W^{I \setminus \{ i \}})_{\af}$
and 
there exists an edge 
$wT_{\xi} \xrightarrow{\ \beta \ } r_{\beta}wT_{\xi}$
in 
$\SiB ^{I \setminus \{ i \}}$
if and only if one of the following conditions holds:
\begin{enumerate}[label=(\Alph*)]\setcounter{enumi}{1}
\item\label{Prop4.11_B} %[(B)]
$c' = c$, 
$\mathsf{T}'(s) = \mathsf{T}(s) + 1$, 
and 
$\mathsf{T}'(u) = \mathsf{T}(u)$ 
for $1 \leq u \leq i$, $u \neq s$.\setcounter{enumi}{16} 
\item\label{Prop4.11_Q} %[(Q)]
$c' = c + 1$, 
$\mathsf{T}'(1) = 1$, 
$\mathsf{T}'(u) = \mathsf{T}(u-1)$ for $2 \leq u \leq i$, 
and 
$\mathsf{T}(i) = n$.
\end{enumerate}
\end{prop}

\begin{rema}
Under the assumptions of Proposition~\ref{prop:tab-cri}, 
the following holds: 
\begin{enumerate}[label=(\arabic*)]
\item\label{Rema4.12_1}
\ref{Prop4.11_B} is equivalent to 
$w(s) \in I$ 
and 
$(\mathsf{T}', c') = f_{w(s)}(\mathsf{T} , c)$
in 
$\mathbb{Y}^{\si}(\varpi_i)$; 
note that 
$\mathsf{T} = \mathsf{T}_w$ 
and 
$\mathsf{T}' = \mathsf{T}_{w r_{\gamma}}$ 
in this case. 
\item\label{Rema4.12_2}
\ref{Prop4.11_Q} is equivalent to 
$(\mathsf{T}', c') = f_0(\mathsf{T} , c)$
in 
$\mathbb{Y}^{\si}(\varpi_i)$; 
note that 
$\mathsf{T} = \mathsf{T}_w$ 
and 
$\mathsf{T}' = \mathsf{T}_{\lfloor w r_{\gamma} \rfloor}$ 
in this case. 
\end{enumerate}
\end{rema}

Proposition~\ref{prop:tab-cri} is established by combining 
Proposition~\ref{prop:phi_J} and Lemmas~\ref{lem:c'=c+chi}--\ref{lem:Q} below.

\begin{lemm} \label{lem:c'=c+chi}
Under the assumptions of Proposition~\ref{prop:tab-cri}, 
we have the following:
\begin{enumerate}[label=(\arabic*)]
\item \label{Lemm_4.13_1}
$\chi = 0$
is equivalent to 
$c' = c$. 
\item\label{Lemm_4.13_2}
$\chi = 1$
is equivalent to 
$c' = c+1$. 
\end{enumerate}
\end{lemm}

\begin{proof}
It suffices to prove that 
$c' = c + \chi$. 
We have 
$r_{\beta} w T_{\xi} 
= r_{w\gamma} t_{\chi w\gamma} w z_{\xi} t_{\xi}
= w r_{\gamma} z_{\xi} t_{\xi + \chi z_{\xi}^{-1} \gamma}$, 
and hence 
$\Pi^{I \setminus \{ i \}}(r_{\beta} w T_{\xi}) 
= \lfloor wr_{\gamma} \rfloor T_{\xi + \chi z_{\xi}^{-1} \gamma}$
by Lemma~\ref{lem:Pi}. 
This gives 
$c' = c_i(\xi + \chi z_{\xi}^{-1} \gamma)
= c_i(\xi) + \chi c_i(z_{\xi}^{-1} \gamma) 
= c + \chi c_i(z_{\xi}^{-1} \gamma)$. 
Since 
$z_{\xi}^{-1} \in W_{I \setminus \{ i \}}$, 
it follows that 
$z_{\xi}^{-1} \gamma \in \Delta^+ \setminus \Delta_{I \setminus \{ i \}}^+$. 
Therefore 
$c_i(z_{\xi}^{-1} \gamma) = 1$
by~\eqref{eq:Delta^i}.
\end{proof}

\begin{lemm}[{\cite[Proposition~2.4.8]{BB}}]
\label{lem:B}
Let 
$w \in W^{I \setminus \{ i \}}$ 
and 
$\gamma = \alpha_{s,t} \in \Delta^+ \setminus \Delta_{I \setminus \{ i \}}^+$, 
$1 \leq s \leq i \leq t \leq n-1$. 
The following conditions are equivalent:
\begin{enumerate}[label=(\arabic*)]
\item\label{Lemm_4.14_1}
$\ell(w r_{\gamma}) = \ell(w) + 1$.
\item\label{Lemm_4.14_2}
$w r_{\gamma} \in W^{I \setminus \{ i \}}$, 
$w r_{\gamma}(s) = w(s) + 1$, 
and 
$w r_{\gamma}(u) = w(u)$ 
for $1 \leq u \leq i$, $u \neq s$.
\end{enumerate}
\end{lemm}

\begin{lemm} \label{lem:Q}
Let 
$w \in W^{I \setminus \{ i \}}$
and 
$\gamma \in \Delta^+ \setminus \Delta_{I \setminus \{ i \}}^+$. 
The following conditions are equivalent:
\begin{enumerate}[label=(\arabic*),series=lemma4.15]
\item\label{lemma4.15_1}
$\ell(\lfloor w r_{\gamma} \rfloor) = \ell(w) + 1 - 2(\gamma,\rho-\rho_{I \setminus \{ i \}})$.
\item\label{lemma4.15_2}
$\lfloor w r_{\gamma} \rfloor (1) = 1$, 
$\lfloor w r_{\gamma} \rfloor (u) = w(u-1)$ for $2 \leq u \leq i$, 
and 
$w(i) = n$.
\end{enumerate}
\end{lemm}

For the proof of Lemma~\ref{lem:Q}, 
we need the following lemma.
Let 
$(i_1 \ i_2 \ \cdots \ i_{l}) \in W$ 
denote the cyclic permutation 
$i_1 \mapsto i_2 \mapsto \cdots \mapsto i_{l} \mapsto i_1$, 
where 
$l \in \BZ_{\geq 2}$ 
and 
$i_1 , i_2 , \ldots , i_{l} \in \{ 1,2,\ldots ,n \}$
are all distinct. 

\goodbreak
\begin{lemm} \label{lem:J^i-ad}\ \\[-1.2em]
\begin{enumerate}[label=(\arabic*)]
\item\label{Lemm_4.16_a}
$Q^{I \setminus \{ i \}} \cap (\Delta^+ \setminus \Delta_{I \setminus \{ i \}}^+) = \{ \alpha_i \}$.
\item\label{Lemm_4.16_b}
$2(\gamma,\rho-\rho_{I \setminus \{ i \}}) = n$ 
for 
$\gamma \in \Delta^+ \setminus \Delta_{I \setminus \{ i \}}^+$.
\item\label{Lemm_4.16_c}
$\ell(w) = (w(1) - 1)i + \sum_{u=2}^i (w(u) - w(u-1) -1)(i - u + 1)$ 
for $w \in W^{I \setminus \{ i \}}$.
\item\label{Lemm_4.16_d}
$z_{\alpha_i}^{I \setminus \{ i \}} = (1 \ 2 \ \cdots \ i)(n \ n-1 \ \cdots \ i+1)$.
\end{enumerate}
\end{lemm}

\begin{proof}
\ref{Lemm_4.16_a}: 
It is clear that 
$\alpha_i \in Q^{I \setminus \{ i \}} \cap (\Delta^+ \setminus \Delta_{I \setminus \{ i \}}^+)$. 
Let 
$\gamma = \alpha_{s,t} \in \Delta^+ \setminus \Delta_{I \setminus \{ i \}}^+$, 
$1 \leq s \leq i \leq t \leq n-1$. 
If $s < i$ (\resp $i < t$), then 
$\alpha_s \in \Delta_{I \setminus \{ i \}}^+$
(\resp $\alpha_t \in \Delta_{I \setminus \{ i \}}^+$) 
and 
$(\gamma , \alpha_s) = 1$
(\resp $(\gamma , \alpha_t) = 1$). 
This proves that 
$\gamma \notin Q^{I \setminus \{ i \}}$ 
unless 
$s = t = i$.

\ref{Lemm_4.16_b}: 
The assertion follows from 
$2(\xi , \rho-\rho_{I \setminus \{ i \}}) = 0$
for 
$\xi \in Q_{I \setminus \{ i \}}$, 
$\gamma \equiv \alpha_i \mod Q_{I \setminus \{ i \}}$
for 
$\gamma \in \Delta^+ \setminus \Delta_{I \setminus \{ i \}}^+$, 
and 
$2(\alpha_i , \rho-\rho_{I \setminus \{ i \}}) = 2(\alpha_i , \rho) - 2(\alpha_i , \rho_{I \setminus \{ i \}}) = 2+ \# (I \setminus \{ i \}) = n$.

\ref{Lemm_4.16_c}: 
This is an immediate consequence of~\eqref{eq:Gr} 
and the fact that the length of 
a permutation equals the number of its inversions
(see~\cite[Proposition~1.5.2]{BB}).

\ref{Lemm_4.16_d}: 
Let 
$I_1 = \{ 1 , \ldots , i-1 \}$
and
$I_2 = \{ i+1 , \ldots , n-1 \}$
be connected components of $I \setminus \{ i \}$. 
We see that
$(i-1,i+1) \in (I_1)_{\af} \times (I_2)_{\af}$ 
satisfies the condition in Lemma~\ref{lem:J-ad} for $\alpha_i$, 
because $\alpha_i \in Q^{I \setminus \{ i \}}$ by~\ref{Lemm_4.16_a}, 
$(\alpha_i , \alpha_{i-1}) = -1$ if $1<i$, 
and 
$(\alpha_i , \alpha_{i+1}) = -1$ if $i<n-1$. 
Therefore 
\begin{align}
z_{\alpha_i}^{I \setminus \{ i \}} 
= 
v_{i-1}^{I_1} v_{i+1}^{I_2}
=
w_0^{I_1} w_0^{I_1 \setminus \{ i-1 \}}
w_0^{I_2} w_0^{I_2 \setminus \{ i+1 \}}.
\end{align}
Now the assertion is shown by the fact that 
the longest element of the symmetric group of degree $N$ 
is the permutation 
$j \mapsto N-j+1$, 
$j \in \{ 1,2,\ldots ,N \}$.
\end{proof}

\begin{proof}[Proof of Lemma~\ref{lem:Q}]
We see from~\cite[Proof of Theorem~10.16]{LS} 
that~\ref{lemma4.15_1} is equivalent~to 
\begin{enumerate}[label=(\arabic*),resume*=lemma4.15]
\item\label{lemma4.15_3} %[(3)]
$\ell(w r_{\gamma}) = \ell(w) + 1 - 2(\gamma,\rho)$
and
$w r_{\gamma}t_{\gamma} \in (W^{I \setminus \{ i \}})_{\af}$.
\end{enumerate}
It follows immediately from Lemmas~\ref{lem:(W^J)_af} 
and~\ref{lem:J^i-ad}~\ref{Lemm_4.16_a} 
that~\ref{lemma4.15_3} is equivalent to 
\begin{enumerate}[label=(\arabic*),resume*=lemma4.15]
\item\label{lemma4.15_4} %[(4)]
$\gamma = \alpha_i$, 
$\ell(w) = \ell(w r_i) + 1$ 
and
$w r_i = \lfloor w r_i \rfloor z_{\alpha_i}^{I \setminus \{ i \}}$.
\end{enumerate}

Let us prove that~\ref{lemma4.15_1} (and~\ref{lemma4.15_4}) imply~\ref{lemma4.15_2}. 
By~\ref{lemma4.15_4} and Lemma~\ref{lem:J^i-ad}~\ref{Lemm_4.16_d}, 
we have
\begin{align} \label{eq:wr_i}
\lfloor w r_i \rfloor = w r_i (z_{\alpha_i}^{I \setminus \{ i \}})^{-1} 
= w(i \ i+1)(i \ \cdots \ 2 \ 1)(i+1 \ \cdots \ n-1 \ n) .
\end{align}
Hence 
$\lfloor w r_i \rfloor (1) = w(i+1)$.

We first assume that $i=1$. 
Then 
$\lfloor w r_1 \rfloor (1) = w(2)$. 
The condition 
$\ell(w) = \ell(w r_1) + 1$ in~\ref{lemma4.15_4}
shows, 
by~\cite[Proposition~1.5.3]{BB}, 
that 
$w(1) > w(2)$. 
Since $w \in W^{I \setminus \{ 1 \}}$, 
it follows from~\eqref{eq:Gr} that 
$w(2) = 1$ 
and, in consequence, 
$\lfloor w r_1 \rfloor (1) = 1$. 
Since 
$\lfloor w r_1 \rfloor \in W^{I \setminus \{ 1 \}}$, 
this implies that 
$\lfloor w r_1 \rfloor = e$ 
and hence 
$w = (n \ n-1 \ \cdots \ 2 \ 1)$ 
by~\eqref{eq:wr_i}. 
This gives 
$w(1) = n$. 

We next assume that 
$1 < i \leq n-1$. 
By~\eqref{eq:wr_i}, 
$\lfloor w r_i \rfloor (u) = w(u-1)$ 
for $2 \leq u \leq i$.
As 
$\lfloor wr_i \rfloor \in W^{I \setminus \{ i \}}$ 
we have
$1 \leq \lfloor wr_i \rfloor (1) < \lfloor wr_i \rfloor (2) = w(1)$. 
Since $w \in W^{I \setminus \{ i \}}$, 
we see from~\eqref{eq:Gr} that 
$w(i+1) = 1$, 
and so 
$\lfloor w r_i \rfloor (1) = 1$. 
It follows from Lemma~\ref{lem:J^i-ad}~~\ref{Lemm_4.16_c} that 
\begin{align}
\ell(\lfloor w r_i \rfloor) 
&= 
(w(1) - 2)(i-1) + \sum_{u=2}^{i-1} \left( w(u) - w(u-1) - 1 \right) (i-u), \\[2mm]
\ell(w) 
&= 
(w(1) - 1)i + \sum_{u=2}^i \left( w(u) - w(u-1) - 1 \right) (i-u+1) ,
\end{align}
which gives 
$\ell(\lfloor w r_i \rfloor) - \ell(w) = 1 - w(i)$. 
By~\ref{lemma4.15_1},~\ref{lemma4.15_4} and Lemma~\ref{lem:J^i-ad}~\ref{Lemm_4.16_b}, 
we have 
$\ell(\lfloor w r_i \rfloor) - \ell(w) = 1 - 2(\alpha_i, \rho-\rho_{I \setminus \{ i \}}) = 1-n$, 
and consequently $w(i) = n$. 

Finally, we prove that~\ref{lemma4.15_2} implies~\ref{lemma4.15_1}. 
In a way similar to the above, 
we have 
$\ell(\lfloor w r_{\gamma} \rfloor) - \ell(w) = 1 - w(i) = 1 - n$. 
Lemma~\ref{lem:J^i-ad}~\ref{Lemm_4.16_b} now shows that~\ref{lemma4.15_1} holds.
\end{proof}

\begin{proof}[Proof of Theorem~\ref{thm:tab-cri}]
If 
$x \preceq y$, 
then 
$c \leq c'$ 
by Proposition~\ref{prop:tab-cri}.
Therefore, we may assume that 
$d := c' - c \geq 0$. 
The proof is by induction on $d$. 

If $d = 0$, then it is obvious from Proposition~\ref{prop:tab-cri} that 
$x \preceq y$ 
is equivalent to 
$\mathsf{T}(u) \leq \mathsf{T}'(u)$
for all 
$1 \leq u \leq i$. 

Let $d > 0$. 
We first assume that 
$\mathsf{T}(u) \leq \mathsf{T}'(u+d)$ 
if 
$1 \leq u \leq i-d$, 
and show that 
$x \preceq y$. 
Let 
$x_1 , x_2 \in (W^{I \setminus \{ i \}})_{\af}$ 
be such that 
$\mathcal{Y}_i(x_1) = (\mathsf{T}_1 , c'), \mathcal{Y}_i(x_2) = (\mathsf{T}_2 , c'-1) \in \mathbb{Y}^{\si}(\varpi_i)$, 
where 
\begin{enumerate}[label=(\arabic*)]
\item
$\mathsf{T}_1(1) = 1$ and 
\begin{align}
\mathsf{T}_1(u) 
=
\begin{cases}
\mathsf{T}'(u) 
&
\text{if} \ 
\mathsf{T}'(u) < n-i+u, \\[2mm]
\mathsf{T}'(u) - 1
&
\text{if} \ 
\mathsf{T}'(u) = n-i+u
\end{cases}
\end{align}
for $2 \leq u \leq i$, 
\item
$\mathsf{T}_2(u) = \mathsf{T}_1(u+1)$ 
for 
$1 \leq u \leq i-1$, 
and 
$\mathsf{T}_2(i) = n$.
\end{enumerate}
By Proposition~\ref{prop:tab-cri}, 
we have 
$x_2 \prec x_1 \preceq y$. 
If we prove that 
\begin{align} \label{eq:T_2>T}
\mathsf{T}(u) \leq \mathsf{T}_2(u+d-1) \ 
\text{if} \ 
1 \leq u \leq i- (d-1), 
\end{align}
then 
$x \preceq x_2$
by induction hypothesis, 
and hence 
$x \preceq y$. 
Note that 
\begin{multline}
\mathsf{T}_2(u+d-1) \\
=
\begin{cases}
\mathsf{T}'(u+d) 
&
\text{if} \ 
1 \leq u+d-1 \leq i-1 \ 
\text{and} \ 
\mathsf{T}'(u+d) < n-i+u+d , \\%[2mm]
\mathsf{T}'(u+d) -1
&
\text{if} \ 
1 \leq u+d-1 \leq i-1 \ 
\text{and} \ 
\mathsf{T}'(u+d) = n-i+u+d , \\%[2mm]
n
&
\text{if} \ 
u+d-1 = i .
\end{cases}
\end{multline}
We prove~\eqref{eq:T_2>T} as follows. 
\begin{enumerate}[label=(\arabic*)]
\item
If
$1 \leq u+d-1 \leq i-1$ 
and 
$\mathsf{T}'(u+d) < n-i+u+d$, 
then 
$\mathsf{T}(u) \leq \mathsf{T}'(u+d) = \mathsf{T}_2(u+d-1)$. 
\item
If
$1 \leq u+d-1 \leq i-1$ 
and 
$\mathsf{T}'(u+d) = n-i+u+d$, 
then 
$\mathsf{T}_2(u+d-1) - \mathsf{T}(u) 
= 
\mathsf{T}'(u+d) - 1 - \mathsf{T}(u)
\geq 
(n-i+u+d) - 1 - (n-i+u)
= d - 1 \geq 0$.
\item
If 
$u+d-1 = i$, 
then 
$\mathsf{T}(u) \leq n = \mathsf{T}_2(u+d-1)$.
\end{enumerate}

We next assume that 
$x \preceq y$, 
and show that
$\mathsf{T}(u) \leq \mathsf{T}'(u+d)$ 
if 
$1 \leq u \leq i-d$. 
We see from Proposition~\ref{prop:tab-cri} that 
there exist 
$x_3 , x_4 \in (W^{I \setminus \{ i \}})_{\af}$ 
such that 
\begin{enumerate}[label=(\arabic*)]
\item
$x \preceq x_4 \prec x_3 \preceq y$, 
\item
$\mathcal{Y}_i(x_3) = (\mathsf{T}_3 , c') , \mathcal{Y}_i(x_4) = (\mathsf{T}_4 , c'-1) \in \mathbb{Y}^{\si}(\varpi_i)$, 
\item
$\mathsf{T}_3(u) \leq \mathsf{T}'(u)$ 
for 
$1 \leq u \leq i$, 
\item
$\mathsf{T}_3(1) = 1$, 
$\mathsf{T}_3(u+1) = \mathsf{T}_4(u)$ 
for 
$1 \leq u \leq i-1$, 
and 
$\mathsf{T}_4(i) = n$. 
\end{enumerate}
By induction hypothesis, 
$\mathsf{T}(u) \leq \mathsf{T}_4(u+d-1)$ 
if 
$1 \leq u \leq i - (d-1)$. 
We have 
$\mathsf{T}'(u+d) - \mathsf{T}(u) \geq \mathsf{T}_3(u+d) - \mathsf{T}_4(u+d-1) = 0$ 
if 
$1 \leq u \leq i-d$. 
\end{proof}

\subsection{Explicit description of crystal structure on \texorpdfstring{$\mathbb{Y}^{\si}(\lambda)$}{Y (infty/2)(lambda)}}

\begin{prop} \label{prop:crystal-str-SiYT}
Let 
$\lambda = \sum_{i \in I} m_i \varpi_i$, 
$m_i \in \BZ_{\geq 0}$, $i \in I$, 
and set 
$N = \sum_{i \in I} m_i$.
Let 
$\mathbb{T} = (\mathsf{T}_1 \mathsf{T}_2 \cdots \mathsf{T}_N , (c_1 , c_2 , \ldots , c_N)) 
\in \CST (\lambda) \times \BZ^N$ 
and 
$j \in I_{\af}$. 
Then 
$\wt (\mathbb{T}) \in P_{\af}$, 
$\varepsilon_j(\mathbb{T}) , \varphi_j(\mathbb{T}) \in \BZ$, and 
$e_j \mathbb{T} , f_j \mathbb{T} \in \CST (\lambda) \times \BZ^N \sqcup \{ \bm{0} \}$
are computed by the following procedure: 
\begin{enumerate}[label=(\roman*)]
\item
$\wt (\mathbb{T}) = \sum_{\nu = 1}^N \wt (\mathsf{T}_{\nu}) - \sum_{\nu = 1}^N c_{\nu} \delta$. 
\item
Let
$\mathsf{T} \in \CST (\varpi_i)$. 
If $j \in I$, then define 
$\epsilon^{(j)}(\mathsf{T}) \in \{ \bm{\oplus}, \bm{\ominus}, \bullet \}$ 
by 
\begin{align}
\epsilon^{(j)}(\mathsf{T})
=
\begin{cases}
\bm{\oplus} & \text{if} \ j \in \mathsf{T} \ \text{and} \ j+1 \notin \mathsf{T}, \\
\bm{\ominus} & \text{if} \ j \notin \mathsf{T} \ \text{and} \ j+1 \in \mathsf{T}, \\
\bullet & \text{otherwise}.
\end{cases}
\end{align}
Likewise, define
$\epsilon^{(0)}(\mathsf{T}) \in \{ \bm{\oplus}, \bm{\ominus}, \bullet \}$ 
by 
\begin{align}
\epsilon^{(0)}(\mathsf{T})
=
\begin{cases}
\bm{\oplus} & \text{if} \ n \in \mathsf{T} \ \text{and} \ 1 \notin \mathsf{T}, \\
\bm{\ominus} & \text{if} \ n \notin \mathsf{T} \ \text{and} \ 1 \in \mathsf{T}, \\
\bullet & \text{otherwise}.
\end{cases}
\end{align}
\item
In 
$\left( \epsilon^{(j)}(\mathsf{T}_1) , \ldots , \epsilon^{(j)}(\mathsf{T}_N) \right)$, 
continue replacing a pair 
$\left( \epsilon^{(j)}(\mathsf{T}_{\nu}) , \epsilon^{(j)}(\mathsf{T}_{\nu'}) \right) = (\bm{\oplus},\bm{\ominus})$ 
with 
$(\bullet , \bullet)$ 
if 
$\nu < \nu'$ 
and 
$\epsilon^{(j)}(\mathsf{T}_{\mu}) = \bullet$ 
for all 
$\nu < \mu < \nu'$ 
until no such pair exists. 
Let 
$\epsilon^{(j)}(\mathbb{T}) \in \{ \bm{\oplus},\bm{\ominus}, \bullet \}^N$
be the resulting tuple such that no 
$\bm{\oplus}$ 
placed to the left of 
$\bm{\ominus}$.
\item
$\varepsilon_j(\mathbb{T})$ 
(\resp $\varphi_j(\mathbb{T})$) 
equals the number of 
$\bm{\ominus}$ (\resp $\bm{\oplus}$) 
in 
$\epsilon^{(j)}(\mathbb{T})$. 
\item
If 
$\bm{\ominus}$ 
is not in 
$\epsilon^{(j)}(\mathbb{T})$, 
then 
$e_j \mathbb{T} = \bm{0}$. 
If there exists 
$\bm{\ominus}$ 
in 
$\epsilon^{(j)}(\mathbb{T})$, 
and the right-most 
$\bm{\ominus}$ 
is at the $\nu$-th place, then
\begin{align}
e_j \mathbb{T} = \left( \mathsf{T}_1 \cdots \mathsf{T}_{\nu -1} \, (e_j \mathsf{T}_{\nu}) \, \mathsf{T}_{\nu +1} \cdots \mathsf{T}_N , 
(c_1 , \ldots , c_{\nu -1} , c_{\nu} - \delta_{j,0} , c_{\nu +1} , \ldots , c_N) \right).
\end{align}
\item
If 
$\bm{\oplus}$ 
is not in 
$\epsilon^{(j)}(\mathbb{T})$, 
then 
$f_j \mathbb{T} = \bm{0}$.
If there exists 
$\bm{\oplus}$ 
in 
$\epsilon^{(j)}(\mathbb{T})$, 
and the left-most 
$\bm{\oplus}$ 
is at the $\nu$-th place, then
\begin{align}
f_j \mathbb{T} = \left( \mathsf{T}_1 \cdots \mathsf{T}_{\nu -1} \, (f_j \mathsf{T}_{\nu}) \, \mathsf{T}_{\nu +1} \cdots \mathsf{T}_N , 
(c_1 , \ldots , c_{\nu -1} , c_{\nu} + \delta_{j,0} , c_{\nu +1} , \ldots , c_N) \right).
\end{align}
\end{enumerate}
\end{prop}

\begin{proof}
Let $\mathsf{T} \in \CST (\varpi_i)$. 
We check at once that the following holds:
\begin{enumerate}[label=(\arabic*)]
\item
$\epsilon^{(j)}(\mathsf{T}) = \bm{\oplus}$ 
if and only if 
$\varepsilon_j(\mathsf{T}) = 0$, 
$\varphi_j(\mathsf{T}) = 1$, and 
$\langle h_j , \wt (\mathsf{T}) \rangle = 1$.
\item
$\epsilon^{(j)}(\mathsf{T}) = \bm{\ominus}$ 
if and only if 
$\varepsilon_j(\mathsf{T}) = 1$, 
$\varphi_j(\mathsf{T}) = 0$, and 
$\langle h_j , \wt (\mathsf{T}) \rangle = -1$.
\item
$\epsilon^{(j)}(\mathsf{T}) = \bullet$ 
if and only if 
$\varepsilon_j(\mathsf{T}) = \varphi_j(\mathsf{T}) = \langle h_j , \wt (\mathsf{T}) \rangle = 0$. 
\end{enumerate}
Then the assertion follows by the same method as in~\cite[\S~2.1]{KN}.
\end{proof}
\begin{thebibliography}{26}
\bibitem{AK} Tatsuya Akasaka and Masaki Kashiwara, Finite-dimensional representations of quantum affine algebras, Publ. Res. Inst. Math. Sci. 33 (1997), no. 5, 839--867.
\bibitem{BN04} Jonathan Beck and Hiraku Nakajima, Crystal bases and two-sided cells of quantum affine algebras, Duke Math. J. 123 (2004), no. 2, 335--402.
\bibitem{BB} Anders Björner and Francesco Brenti, Combinatorics of Coxeter groups, Graduate Texts in Mathematics, vol. 231, Springer, New York, 2005.
\bibitem{F} William Fulton, Young tableaux. with applications to representation theory and geometry, London Mathematical Society Student Texts, vol. 35, Cambridge University Press, Cambridge, 1997.
\bibitem{I} Motohiro Ishii, Path model for representations of generalized Kac--Moody algebras, J. Algebra 379 (2013), 277--300.
\bibitem{INS16} Motohiro Ishii, Satoshi Naito, and Daisuke Sagaki, Semi-infinite Lakshmibai--Seshadri path model for level-zero extremal weight modules over quantum affine algebras, Adv. Math. 290 (2016), 967--1009.
\bibitem{J} Anthony Joseph, Quantum groups and their primitive ideals, Ergebnisse der Mathematik und ihrer Grenzgebiete (3) [Results in Mathematics and Related Areas (3)], vol. 29, Springer-Verlag, Berlin, 1995.
\bibitem{JL} Anthony Joseph and Polyxeni Lamprou, A Littelmann path model for crystals of generalized Kac--Moody algebras, Adv. Math. 221 (2009), no. 6, 2019--2058.
\bibitem{Kac} Victor G. Kac, Infinite dimensional Lie algebras, Results in Mathematics and Related Areas (3), Cambridge University Press, Cambridge, 1990, 3rd edition.
\bibitem{KMN} Seok-Jin Kang, Masaki Kashiwara, Kailash C. Misra, Tetsuji Miwa, Toshiki Nakashima, and Atsushi Nakayashiki, Perfect crystals of quantum affine Lie algebras, Duke Math. J. 68 (1992), no. 3, 499--607.
\bibitem{K94} Masaki Kashiwara, Crystal bases of modified quantized enveloping algebra, Duke Math. J. 73 (1994), no. 2, 383--413.
\bibitem{K95} Masaki Kashiwara, On crystal bases, in Representations of groups (Banff, AB, 1994), CMS Conf. Proc., vol. 16, Amer. Math. Soc., Providence, RI, 1995, pp. 155--197.
\bibitem{K96} Masaki Kashiwara, Similarity of crystal bases, in Lie algebras and their representations (Seoul, 1995), Contemp. Math., vol. 194, Amer. Math. Soc., Providence, RI, 1996, pp. 177--186.
\bibitem{K02} Masaki Kashiwara, On level-zero representations of quantized affine algebras, Duke Math. J. 112 (2002), no. 1, 117--175.
\bibitem{KN} Masaki Kashiwara and Toshiki Nakashima, Crystal graphs for representations of the $q$-analogue of classical Lie algebras, J. Algebra 165 (1994), no. 2, 295--345.
\bibitem{KNS} Syu Kato, Satoshi Naito, and Daisuke Sagaki, Equivariant K-theory of semi-infinite flag manifolds and Pieri--Chevalley formula, \url{https://arxiv.org/abs/1702.02408}, 2017.
\bibitem{Kw} Jae-Hoon Kwon, Crystal graphs and the combinatorics of Young tableaux, in Handbook of Algebra, vol. 6, Elsevier/North-Holland, Amsterdam, 2009, pp. 473--504.
\bibitem{LS} Thomas Lam and Mark Shimozono, Quantum cohomology of $G/P$ and homology of affine Grassmannian, Acta Math. 204 (2010), no. 1, 49--90.
\bibitem{LNSSS15} Cristian Lenart, Satoshi Naito, Daisuke Sagaki, Anne Schilling, and Mark Shimozono, A uniform model for Kirillov--Reshetikhin crystals I: Lifting the parabolic quantum Bruhat graph, Int. Math. Res. Not. IMRN (2015), no. 7, 1848--1901.
\bibitem{LNSSS16} Cristian Lenart, Satoshi Naito, Daisuke Sagaki, Anne Schilling, and Mark Shimozono, Quantum Lakshmibai--Seshadri paths and root operators, in Schubert calculus---Osaka 2012, Adv. Stud. Pure Math., vol. 71, Math. Soc. Japan, Tokyo, 2016, pp. 267--294.
\bibitem{L94} Peter Littelmann, A Littlewood--Richardson rule for symmetrizable Kac--Moody algebras, Invent. Math. 116 (1994), no. 1-3, 329--346.
\bibitem{L95} Peter Littelmann, Paths and root operators in representation theory, Ann. of Math. (2) 142 (1995), no. 3, 499--525.
\bibitem{L96} Peter Littelmann, A plactic algebra for semisimple Lie algebras, Adv. Math. 124 (1996), no. 2, 312--331.
\bibitem{N} Hiraku Nakajima, Extremal weight modules of quantum affine algebras, in Representation theory of algebraic groups and quantum groups, Adv. Stud. Pure Math., vol. 40, Math. Soc. Japan, Tokyo, 2004, pp. 343--369.
\bibitem{P} Dale Peterson, Quantum cohomology of $G/P$, Lect. Notes, Massachusetts Institute of Technology, Cambridge, MA, 1997.
\bibitem{S} Mark Shimozono, Affine type $A$ crystal structure on tensor products of rectangles, Demazure characters, and nilpotent varieties, J. Algebraic Combin. 15 (2002), no. 2, 151--187.
\end{thebibliography}
\end{document}
